\documentclass[UTF8, 12pt, fontset=fandol, oneside]{ctexart}
\usepackage{researchnotes}

% 保留分册独立编译时的章号前缀。
\renewcommand{\thesection}{7.\arabic{section}}

\title{第七章\quad 定积分}
\author{}
\date{}

\begin{document}

\maketitle

\section{定积分的概念与微积分基本定理}

\subsection{曲边梯形的面积}

定积分最朴素的思想可以追溯到阿基米德时代, 阿基米德的 ``由直线组成平面图形, 由平面组成立体图形'' 的思想可以认为是定积分的萌芽. 他曾借助圆柱和圆锥的体积并利用细分的方法求出了球的体积. 但由于用细分的方法无法解决一般函数的定积分计算问题, 因此定积分的研究在漫长的岁月里一直没有进展. 直到 17 世纪, 在天文学和物理学等学科的促进下, 数学研究才取得了很大的进展. 伽利略和开普勒的一系列发现, 如开普勒的行星运动三定律等, 导致了导数的诞生. 导数对现代数学理论的建立起了巨大的促进作用. 到了 17 世纪下半叶, 牛顿和莱布尼茨两人在综合、发展前人工作的基础上, 几乎同时建立了微积分理论. 正是这种理论, 使得定积分有了简便的计算方法, 从而也使得它具有极大的理论和实用价值。

可以引入定积分概念的问题很多, 如开普勒研究的椭圆面积问题; 已知速度求路程的问题; 已知密度求质量的问题; 等等. 下面我们以求由区间 $[a,b]$ 上非负函数 $f(x)$ 的图像(曲线)与 $x$ 轴及直线 $x=a,x=b$ 所围图形 $S$(如图 7.1.1, 我们称该图形为曲边梯形)的面积为例来详细阐述定积分的基本思想。

回忆一下, 我们目前所能精确求出面积的图形只能是由有限条线段所围成的平面图形. 如果 $y=f(x)$ 不是一个线性函数, 则不可能简单地用某个我们已知的计算公式算出该曲边梯形的面积. 因此我们需要探究如何来求这种图形面积的方法。

为了使问题简化, 我们假定 $f(x)$ 在区间 $[a,b]$ 上连续, 并记 $M,m$ 分别为它在区间 $[a,b]$ 上的最大值和最小值. 在重积分理论中, 我们将讨论一般的平面图形的面积问题并证明曲边梯形 $S$ 的面积必定存在. 现在我们假定该曲边梯形的面积存在并将其面积仍记为 $S$, 易于看出
\[
  m(b-a)\leq S\leq M(b-a).
\]

设函数 $y=f(x)(x\in[a,b])$ 的图像为 $T$. 由连续函数的性质可以知道 $T$ 上必有一点 $(c,f(c))$, 使得 $S=f(c)(b-a)$. 我们事先无法知道 $c$ 在何处, 能做的是在区间 $[a,b]$ 上取一点 $\xi$, 计算 $f(\xi)(b-a)$. $f(\xi)(b-a)$ 的几何意义很明显, 它是以区间 $[a,b]$ 为底, 以 $f(\xi)$ 为高的矩形的面积. 它与 $S$ 的误差可以由下式估计:
\[
  |S-f(\xi)(b-a)|=|f(c)-f(\xi)|(b-a)\leq (M-m)(b-a),
\]
其中 $M-m$ 也称为 $f(x)$ 在 $[a,b]$ 上的振幅。

如果我们等分区间 $[a,b]$ 成两个小区间, 并令 $x_1=\dfrac{a+b}{2}$(见图 7.1.1), 则直线 $x=x_1$ 将 $S$ 分成两个曲边梯形 $S_1$ 及 $S_2$. 现将它们的面积仍然分别记为 $S_1$ 和 $S_2$, 显然有
\[
  S=S_1+S_2.
\]
在 $S_1$ 和 $S_2$ 中重复我们刚才在区间 $[a,b]$ 上所做的过程, 记 $M_1,M_2$ 分别为 $f(x)$ 在区间 $[a,x_1]$ 和 $[x_1,b]$ 上的最大值, $m_1,m_2$ 分别为 $f(x)$ 在区间 $[a,x_1]$ 和 $[x_1,b]$ 上的最小值. 在区间 $[a,x_1]$ 上任意取 $\xi_1$, 在区间 $[x_1,b]$ 上任意取 $\xi_2$, 则有以下的估计式:
\[
\begin{aligned}
  |f(\xi_1)(x_1-a)-S_1|&\leq (M_1-m_1)\frac{b-a}{2},\\
  |f(\xi_2)(b-x_1)-S_2|&\leq (M_2-m_2)\frac{b-a}{2}.
\end{aligned}
\]
因此
\[
\begin{aligned}
  &|f(\xi_1)(x_1-a)+f(\xi_2)(b-x_1)-S|\\
  &\quad\leq [(M_1-m_1)+(M_2-m_2)]\frac{b-a}{2}\\
  &\quad\leq (M-m)(b-a).
\end{aligned}
\]

上述成立是因为分割后的误差估计由小区间上的振幅给出, 而一个函数在小区间上的振幅不会比在包含该区间的大区间上的振幅大的缘故. 上式还告诉我们: 对于分成两块所得到的和式 $f(\xi_1)(x_1-a)+f(\xi_2)(b-x_1)$ 与 $S$ 的误差和原来 $f(\xi)(b-a)$ 与 $S$ 的误差, 分割后的误差的估计式不会变大. 在许多情况下, 比如 $f(x)$ 在区间 $[a,b]$ 单调, 分割后任意取值得到的和式与 $S$ 的误差估计比分割前的误差估计会变得更小。

当区间 $[a,b]$ 分割后的区间长度很小时, 由 $f(x)$ 的一致连续性, $f(x)$ 在每个小区间上的振幅都可以很小. 基于这种思想, 我们将区间 $[a,b]$ 分割成更小的区间, 如对充分大的 $n$, 将区间 $[a,b]$ 等分为 $n$ 个小区间(见图 7.1.2). 设其分点为
\[
  a=x_0<x_1<x_2<\cdots<x_n=b,
\]
记 $f(x)$ 在区间 $[x_{i-1},x_i](i=1,2,\cdots,n)$ 上的最大值为 $M_i$, 最小值为 $m_i$. 在每个区间上任取 $\xi_i$, 作和式
\[
  \sum_{i=1}^n f(\xi_i)(x_i-x_{i-1});
\]
则此和式与 $S$ 的误差可以有下面的估计:
\begin{equation}
  \left|\sum_{i=1}^n f(\xi_i)(x_i-x_{i-1})-S\right|
  \leq \sum_{i=1}^n (M_i-m_i)(x_i-x_{i-1}).
  \tag{7.1.1}
\end{equation}
上述不等式右边即为图 7.1.2 中与曲线相交的小矩形面积之和. 由于 $f(x)$ 在区间 $[a,b]$ 上连续, 从而一致连续, 因此对于 $\forall\varepsilon>0$, $\exists\delta>0$, 当每个 $x_i-x_{i-1}<\delta$ 时(这只要 $n$ 充分大即可), 有 $M_i-m_i<\varepsilon$. 于是我们有
\[
  \left|\sum_{i=1}^n f(\xi_i)(x_i-x_{i-1})-S\right|
  \leq \sum_{i=1}^n (M_i-m_i)(x_i-x_{i-1})<\varepsilon(b-a).
\]

以上分析告诉我们, 只要 $f(x)$ 在区间 $[a,b]$ 上连续, 当区间 $[a,b]$ 分割得充分 ``细'' 时, 对每个小曲边梯形的曲边 ``以直代曲'', 作出的和式的值可以任意地接近该曲边梯形的面积. 因此有
\[
  \lim_{n\to\infty}\sum_{i=1}^n f(\xi_i)(x_i-x_{i-1})=S.
\]

下面我们以由曲线 $f(x)=x^2(x\in[0,1])$ 与 $x$ 轴及直线 $x=1$ 所围的图形为例来求其面积 $S$(见图 7.1.3)。

将区间 $[0,1]$ 进行 $n$ 等分, 其分点为 $0,\dfrac1n,\dfrac2n,\cdots,\dfrac{n-1}{n},1$. 在区间 $\left[\dfrac{i-1}{n},\dfrac{i}{n}\right](i=1,2,\cdots,n)$ 上任取 $\xi_i$, 作和式
\[
  \sum_{i=1}^n f(\xi_i)\left(\frac{i}{n}-\frac{i-1}{n}\right)
  =\frac1n\sum_{i=1}^n \xi_i^2.
\]
由于 $\dfrac{i-1}{n}\leq \xi_i\leq \dfrac{i}{n}(i=1,2,\cdots,n)$, 且 $f(x)$ 在区间 $[0,1]$ 单调上升, 因此有
\[
  m_i=\left(\frac{i-1}{n}\right)^2\leq \xi_i^2\leq \left(\frac{i}{n}\right)^2=M_i.
\]
注意到
\[
  \lim_{n\to\infty}\frac1{n^3}\sum_{i=1}^n (i-1)^2
  =\lim_{n\to\infty}\frac1{n^3}\sum_{i=1}^n i^2
  =\frac13,
\]
我们有
\[
  \left|\frac1n\sum_{i=1}^n \xi_i^2-S\right|
  \leq \left|\frac1{n^3}\sum_{i=1}^n i^2-\frac1{n^3}\sum_{i=1}^n (i-1)^2\right|
  \to 0\quad(n\to\infty).
\]
因此有
\[
  S=\lim_{n\to\infty}\frac1n\sum_{i=1}^n \xi_i^2=\frac13.
\]

我们再来举一个质点做直线运动的例子. 假定一个质点在直线上朝着一个方向运动, 其速度 $v=v(t)$ 是时间 $t$ 的函数. 为简便计, 假定 $v(t)\geq 0$. 现在我们来求在时间段 $[a,b]$ 内质点走过的路程 $s$。

若 $v(t)=c$ 是常数函数, 则 $s=c(b-a)$. 如果 $v(t)$ 不是常数函数, 即质点做非匀速运动时, 我们如何来求 $s$ 呢? 为此我们来分析一下, 假定 $v(t)$ 是连续变化的(在许多实际问题中大都是如此), 则对任意 $t_0\in[a,b]$, $v(t)$ 在 $t_0$ 附近的值与 $v(t_0)$ 很接近. 换句话说, 对很小的 $\delta>0$, 在时间段 $[t_0-\delta,t_0+\delta]$ 内, 质点走过的路程 $\tilde s\approx 2\delta v(t_0)$(简称 ``以不变代变''). 这就启发我们将区间 $[a,b]$ 作分割:
\[
  a=t_0<t_1<\cdots<t_n=b,
\]
当 $\lambda=\max\limits_{1\leq i\leq n}\{t_i-t_{i-1}\}$ 很小时, 在每个时间段 $[t_{i-1},t_i](i=1,2,\cdots,n)$ 内质点走过的路程 $s_i\approx v(\xi_i)(t_i-t_{i-1})$, 其中 $\xi_i\in[t_{i-1},t_i]$. 因此, 总的路程
\begin{equation}
  s=\sum_{i=1}^n s_i\approx \sum_{i=1}^n v(\xi_i)(t_i-t_{i-1}).
  \tag{7.1.2}
\end{equation}
当时间段 $[a,b]$ 的分割对应的小时间段的长度趋于零时, 上面路程的近似值与路程的准确值无限地接近. 因此若和式 (7.1.2) 有极限, 此极限便是我们要求的路程. 容易看出以上已知速度函数求路程与我们在前面求曲边梯形的面积具有相同的 ``纯数学'' 极限过程。

\subsection{定积分的定义}

上面我们对若干问题进行了讨论, 在自然界还有许多类似的问题, 在解决它们的过程中可以抽象出一个相同的 ``纯数学'' 极限过程. 为此我们给出下面的定义。

\noindent\textbf{定义 7.1.1}\quad 设函数 $f(x)$ 在区间 $[a,b]$ 上有定义, 对于区间 $[a,b]$ 的一个分割
\[
  \Delta:a=x_0<x_1<\cdots<x_n=b,
\]
记 $\Delta x_i=x_i-x_{i-1}(i=1,2,\cdots,n)$, $\lambda(\Delta)=\max\limits_{1\leq i\leq n}\{\Delta x_i\}$. 在每个小区间 $[x_{i-1},x_i](i=1,2,\cdots,n)$ 上任取 $\xi_i$, 作和式
\[
  \sum_{i=1}^n f(\xi_i)\Delta x_i.
\]
如果当 $\lambda(\Delta)\to0$ 时, 上述和式存在极限 $I$, 且 $I$ 不依赖分割 $\Delta$ 的选取及 $\xi_i(i=1,2,\cdots,n)$ 在 $[x_{i-1},x_i]$ 上的选取, 则称 $f(x)$ 在区间 $[a,b]$ 上是黎曼可积(简称可积)的, 同时称 $I$ 为 $f(x)$ 在区间 $[a,b]$ 上的定积分, 记为
\[
  I=\int_a^b f(x)\mathrm dx,
\]
其中 $a$ 与 $b$ 分别称为定积分的下限和上限, $f(x)$ 称为被积函数, $x$ 称为积分变量。

在定积分理论的讨论中, 我们经常要遇到分割求和. 为了使行文简洁, 在今后对每个给定的区间 $[a,b]$, 我们总是用
\[
  \Delta:a=x_0<x_1<\cdots<x_n=b
\]
来表示它的一个分割, 对于这个分割总是记 $\Delta x_i=x_i-x_{i-1}(i=1,2,\cdots,n)$, 并且记 $\lambda(\Delta)=\max\limits_{1\leq i\leq n}\{\Delta x_i\}$. 在 $[x_{i-1},x_i](i=1,2,\cdots,n)$ 上任取一点 $\xi_i$, 我们称和式 $\sum\limits_{i=1}^n f(\xi_i)\Delta x_i$ 为函数 $f(x)$ 关于 $\Delta$ 的黎曼和。

用 ``$\varepsilon$-$\delta$'' 语言叙述定积分的定义则更为精确。

\noindent\textbf{定义 7.1.1'}\quad 设函数 $f(x)$ 在区间 $[a,b]$ 上有定义, 若存在常数 $I\in\mathbb R$, 使得对于 $\forall\varepsilon>0$, $\exists\delta>0$, 对区间 $[a,b]$ 的任何一个分割 $\Delta$, 当 $\lambda(\Delta)<\delta$ 时, 在每个 $[x_{i-1},x_i](i=1,2,\cdots,n)$ 上任取 $\xi_i$, 都有
\[
  \left|\sum_{i=1}^n f(\xi_i)\Delta x_i-I\right|<\varepsilon,
\]
则称函数 $f(x)$ 在区间 $[a,b]$ 是黎曼可积的, 并称 $I$ 为函数 $f(x)$ 在区间 $[a,b]$ 上的定积分。

今后, 我们用记号 $f(x)\in R[a,b]$ 表示函数 $f(x)$ 在区间 $[a,b]$ 上可积。

从定义 7.1.1(或定义 7.1.1') 可以看出: 若函数 $f(x)$ 在区间 $[a,b]$ 上可积, 则它在区间 $[a,b]$ 上的定积分 $\int_a^b f(x)\mathrm dx$ 是一种特殊和式的极限. 因此, 和式写成 $\sum\limits_{i=1}^n f(\xi_i)\Delta x_i$ 与 $\sum\limits_{i=1}^n f(\tau_i)\Delta t_i$ 并不影响和式的极限, 此和式是由函数关系 $f$ 决定的, 与自变量取 $x$ 还是取 $t$ 无关, 即
\[
  \int_a^b f(x)\mathrm dx=\int_a^b f(t)\mathrm dt.
\]

值得注意的是, 上述和式的极限与我们所学过的序列极限和函数极限是不同的. 在 $\lambda(\Delta)\to0$ 的过程中, 这个极限过程具有两个任意性, 即分割的任意性及在每个小区间上 $\xi_i$ 选取的任意性. 尽管这些极限过程有所不同, 但定积分定义中和式的极限与序列极限或函数极限仍具有一些相似的性质, 如一个函数的定积分存在, 则它必定唯一. 这些性质的证明可以由定义直接推出, 读者可自证之。

\subsection{定积分的几何意义}

设函数 $f(x)$ 在区间 $[a,b]$ 上连续. 我们先看 $f(x)\geq0$ 的情形. 根据我们前面的讨论可知: 定积分 $\int_a^b f(x)\mathrm dx$ 是由直线 $x=a,x=b,y=0$ 与曲线 $y=f(x)$ 所围成的曲边梯形的面积 $S$。

当 $f(x)\leq0$ 时, 同样地设由直线 $x=a,x=b,y=0$ 与曲线 $y=f(x)$ 所围成的曲边梯形的面积为 $S$, 则此时
\[
  S=-\int_a^b f(x)\mathrm dx.
\]
当 $f(x)$ 有正有负时, 则 $\int_a^b f(x)\mathrm dx$ 是由直线 $x=a,x=b,y=0$ 与 $y=f(x)$ 所围成的几个曲边梯形中, 位于 $x$ 轴上方的各个曲边梯形面积之和减去位于 $x$ 轴下方的各个曲边梯形面积之和. 如图 7.1.4, 则有
\[
  \int_a^b f(x)\mathrm dx=S_1-S_2+S_3-S_4.
\]

\subsection{连续函数的可积性}

在对曲边梯形面积的讨论中, 如果假定该面积存在, 我们找到了一种求出该面积的方法. 有了这一基础, 对一个连续函数 $y=f(x),x\in[a,b]$, 我们比较容易证明它在区间 $[a,b]$ 上必定是可积的。

\noindent\textbf{定理 7.1.1}\quad 设函数 $f(x)\in C[a,b]$, 则 $f(x)\in R[a,b]$。

\noindent\textbf{证明}\quad 记 $m,M$ 分别为函数 $f(x)$ 在区间 $[a,b]$ 上的最小和最大值. 我们注意到, 对区间 $[a,b]$ 的任意分割 $\Delta$, $f(x)$ 关于 $\Delta$ 的任一黎曼和满足
\begin{equation}
  m(b-a)\leq \sum_{i=1}^n m_i\Delta x_i
  \leq \sum_{i=1}^n f(\xi_i)\Delta x_i
  \leq \sum_{i=1}^n M_i\Delta x_i
  \leq M(b-a),
  \tag{7.1.3}
\end{equation}
其中 $m_i$ 与 $M_i$ 分别为 $f(x)$ 在区间 $[x_{i-1},x_i](i=1,2,\cdots,n)$ 上的最小值和最大值. 特别地, 对每一个 $n\in\mathbb N$, 将区间 $[a,b]$ 的 $n$ 等分的分割记为 $\Delta_n$, 在此分割的每个小区间 $[x_{j-1},x_j](j=1,2,\cdots,n)$ 上取定 $\xi_j=x_j$, 我们可得一个特殊的黎曼和
\[
  S_n=\sum_{j=1}^n f(x_j)\frac{b-a}{n}.
\]
这样一来, 我们便得到了一个序列 $\{S_n\}$, 由 (7.1.3) 知它是一个有界序列, 因此存在收敛子列 $\{S_{n_k}\}\subset\{S_n\}$. 设 $\lim\limits_{k\to\infty}S_{n_k}=S$, 下面我们证明 $S$ 即为函数 $f(x)$ 在区间 $[a,b]$ 上的定积分。

对于 $\forall\varepsilon>0$, 由函数 $f(x)$ 在区间 $[a,b]$ 上一致连续知, $\exists\delta>0$, 当 $x',x''\in[a,b]$ 且 $|x'-x''|<\delta$ 时, 有
\[
  |f(x')-f(x'')|<\frac{\varepsilon}{4(b-a)},
\]
从而对任意的分割 $\Delta$, 当 $\lambda(\Delta)<\delta$ 时, 有
\[
  \sum_{i=1}^n (M_i-m_i)\Delta x_i<\frac{\varepsilon}{4}.
\]
由于 $\lim\limits_{k\to\infty}S_{n_k}=S$, 因此存在 $K>0$, 使得当 $n_k>K$ 时, 有 $\lambda(\Delta_{n_k})<\delta$ 和 $|S_{n_k}-S|<\dfrac{\varepsilon}{2}$ 同时成立. 现取定一个 $n_{k_0}>K$, 使得它同时满足这两个条件, 即 $\lambda(\Delta_{n_{k_0}})<\delta$ 和 $|S_{n_{k_0}}-S|<\dfrac{\varepsilon}{2}$。

对于任意的分割 $\Delta$, 当 $\lambda(\Delta)<\delta$ 时, 我们有以下的估计:
\begin{equation}
\begin{aligned}
  \left|\sum_{i=1}^n f(\xi_i)\Delta x_i-S\right|
  &\leq \left|\sum_{i=1}^n f(\xi_i)\Delta x_i-S_{n_{k_0}}\right|+|S_{n_{k_0}}-S|\\
  &< \left|\sum_{i=1}^n f(\xi_i)\Delta x_i-\sum_{j=1}^{n_{k_0}} f(x_j)\frac{b-a}{n_{k_0}}\right|+\frac{\varepsilon}{2}.
\end{aligned}
\tag{7.1.4}
\end{equation}
为了估计上述不等式右边两个黎曼和的差, 记 $\Delta'$ 为 $\Delta$ 与 $\Delta_{n_{k_0}}$ 的分点集的并所构成的区间 $[a,b]$ 的分割, 显然有 $\lambda(\Delta')<\delta$. 容易看出对关于 $\Delta',\Delta,\Delta_{n_{k_0}}$ 的任意黎曼和, 有下述估计:
\[
  \left|\sum_{i=1}^n f(\xi_i)\Delta x_i-\sum_{k=1}^{n'}f(\xi_k')\Delta x_k'\right|
  \leq \sum_{i=1}^n(M_i-m_i)\Delta x_i<\frac{\varepsilon}{4}
\]
和
\[
  \left|\sum_{j=1}^{n_{k_0}}f(x_j)\frac{b-a}{n_{k_0}}-\sum_{k=1}^{n'}f(\xi_k')\Delta x_k'\right|
  \leq \sum_{j=1}^{n_{k_0}}(M_j'-m_j')\frac{b-a}{n_{k_0}}<\frac{\varepsilon}{4},
\]
其中 $\sum\limits_{k=1}^{n'}f(\xi_k')\Delta x_k'$ 是关于 $\Delta'$ 的任意一个黎曼和, $\sum\limits_{i=1}^n(M_i-m_i)\Delta x_i$ 是关于分割 $\Delta$ 求和, $\sum\limits_{j=1}^{n_{k_0}}(M_j'-m_j')\dfrac{b-a}{n_{k_0}}$ 是关于分割 $\Delta_{n_{k_0}}$ 求和. 因此
\[
\begin{aligned}
  \left|\sum_{i=1}^n f(\xi_i)\Delta x_i-\sum_{j=1}^{n_{k_0}}f(x_j)\frac{b-a}{n_{k_0}}\right|
  &\leq \left|\sum_{i=1}^n f(\xi_i)\Delta x_i-\sum_{k=1}^{n'}f(\xi_k')\Delta x_k'\right|\\
  &\quad+\left|\sum_{j=1}^{n_{k_0}}f(x_j)\frac{b-a}{n_{k_0}}-\sum_{k=1}^{n'}f(\xi_k')\Delta x_k'\right|\\
  &<\frac{\varepsilon}{4}+\frac{\varepsilon}{4}=\frac{\varepsilon}{2}.
\end{aligned}
\]
将此估计代入 (7.1.4) 式即得
\[
  \left|\sum_{i=1}^n f(\xi_i)\Delta x_i-S\right|<\varepsilon.
\]
这就证明了函数 $f(x)$ 在区间 $[a,b]$ 上可积, 并且 $\int_a^b f(x)\mathrm dx=S$. 证毕。

\subsection{微积分基本定理}

从定积分的定义可知, 定积分应在许多问题中有重要应用, 因此如何计算一个函数的定积分是一个十分关键的问题. 如果对每个函数都要从定积分的定义出发, 对区间分割, 在分割后的每个小区间任取一点作和式, 再求极限, 当函数稍微复杂一点, 求定积分的值就会面临巨大的困难. 如果没有一种相对简单的定积分计算方法, 定积分或许不会成为一种有用的理论和工具。

如何发现定积分简便的计算方法呢? 我们可以先考查一下以下具体问题. 假定我们已知一个质点做直线运动, 其速度函数为 $v(t)>0(t\in[t_0,t_1])$, 现在来求它在该时间段内走过的路程. 前面我们已经知道, 若速度函数 $v(t)$ 的性质比较好(如 $v(t)$ 是连续的), $v(t)$ 在 $[t_0,t_1]$ 上的定积分
\[
  \int_{t_0}^{t_1}v(t)\mathrm dt
\]
刚好是质点在该时间段内走过的路程. 从另外一方面来看, 该段路程刚好是路程函数 $s(t)$ 在 $t=t_1$ 和 $t=t_0$ 两点的值之差. 换句话说, 我们有
\[
  \int_{t_0}^{t_1}v(t)\mathrm dt=s(t_1)-s(t_0).
\]
在微分学理论中, 我们已经知道了 $s'(t)=v(t)$。

那么, 若一个函数 $f(x)$ 在区间 $[a,b]$ 上可积, 且存在原函数 $F(x)$, 是否总是成立
\[
  \int_a^b f(x)\mathrm dx=F(b)-F(a)
\]
呢? 这个问题已由牛顿和莱布尼茨解决. 因此人们称下面定理给出的计算定积分的公式为牛顿-莱布尼茨公式. 由于它是微积分理论中最重要的结果, 人们也称它为微积分基本定理。

\noindent\textbf{定理 7.1.2(微积分基本定理)}\quad 设函数 $f(x)$ 在区间 $[a,b]$ 上有定义, 并且满足以下两个条件:

(1) 在区间 $[a,b]$ 上可积;

(2) 在区间 $[a,b]$ 上存在原函数 $F(x)$,

则有
\begin{equation}
  \int_a^b f(x)\mathrm dx=F(b)-F(a)\triangleq F(x)\big|_a^b.
  \tag{7.1.5}
\end{equation}

\noindent\textbf{证明}\quad 首先, 由于函数 $f(x)$ 在区间 $[a,b]$ 上可积, 任取区间 $[a,b]$ 的分割
\[
  \Delta:a=x_0<x_1<\cdots<x_n=b,
\]
在 $[x_{i-1},x_i](i=1,2,\cdots,n)$ 上任取一点 $\xi_i$, 则有
\[
  \lim_{\lambda(\Delta)\to0}\sum_{i=1}^n f(\xi_i)\Delta x_i=\int_a^b f(x)\mathrm dx.
\]
其次, 对于分割 $\Delta$, 我们有
\[
  F(b)-F(a)=\sum_{i=1}^n[F(x_i)-F(x_{i-1})].
\]
在 $[x_{i-1},x_i](i=1,2,\cdots,n)$ 上对 $F(x)$ 应用拉格朗日微分中值定理得
\[
  F(x_i)-F(x_{i-1})=f(\xi_i')\Delta x_i,
\]
其中 $\xi_i'\in(x_{i-1},x_i)$. 因此我们有
\[
\begin{aligned}
  F(b)-F(a)
  &=\lim_{\lambda(\Delta)\to0}\sum_{i=1}^n[F(x_i)-F(x_{i-1})]\\
  &=\lim_{\lambda(\Delta)\to0}\sum_{i=1}^n f(\xi_i')\Delta x_i\\
  &=\int_a^b f(x)\mathrm dx.
\end{aligned}
\]
证毕。

利用牛顿-莱布尼茨公式可以使许多定积分的计算变得十分简单。

例如前面我们利用分割、求和以及计算极限求出了 $\int_0^1x^2\mathrm dx=\dfrac13$. 由于 $f(x)=x^2\in C[0,1]$, 由定理 7.1.1, 我们知 $f(x)=x^2\in R[0,1]$. 在不定积分中, 我们知道 $F(x)=\dfrac{x^3}{3}$ 是函数 $f(x)$ 的一个原函数, 现利用牛顿-莱布尼茨公式即有
\[
  \int_0^1 x^2\mathrm dx=\left.\frac{x^3}{3}\right|_0^1=\frac13.
\]

\noindent\textbf{例 7.1.1}\quad 设
\[
  I_n=\frac1{n^{\alpha+1}}(1+2^\alpha+3^\alpha+\cdots+n^\alpha)\quad(\alpha>0),
\]
求 $\lim\limits_{n\to\infty}I_n$。

\noindent\textbf{解}\quad 由于
\[
  I_n=\frac1n\left(\frac1{n^\alpha}+\frac{2^\alpha}{n^\alpha}+\frac{3^\alpha}{n^\alpha}+\cdots+\frac{n^\alpha}{n^\alpha}\right)
  =\frac1n\sum_{i=1}^n\frac{i^\alpha}{n^\alpha},
\]
因此 $I_n$ 可以看成是函数 $x^\alpha$ 在区间 $[0,1]$ 上关于分割
\[
  0<\frac1n<\frac2n<\cdots<\frac nn=1
\]
的一个黎曼和. 由 $x^\alpha$ 在 $[0,1]$ 上连续, 知它在 $[0,1]$ 上可积. 在不定积分中, 我们已知 $\dfrac1{\alpha+1}x^{\alpha+1}$ 是 $x^\alpha$ 的一个原函数, 因此由牛顿-莱布尼茨公式得
\[
  \lim_{n\to\infty}I_n=\int_0^1 x^\alpha\mathrm dx
  =\left.\frac1{\alpha+1}x^{\alpha+1}\right|_0^1
  =\frac1{\alpha+1}.
\]

\section{可积性问题}

我们已经指出, 在微积分的理论中, 微积分基本定理是最重要的定理. 从该定理可以看出, 在定积分理论中我们要解决的两个基本问题是:

(1) 函数的可积性;

(2) 原函数的存在性。

在以下各节中, 我们将逐步展开来讨论这两个问题. 在本节中, 我们将讨论函数的可积性问题. 尽管我们已经证明了连续函数的可积性, 但在许多数学理论中我们还需要讨论更为广泛的可积函数类. 事实上, 正是在研究更为广泛的函数的可积性时, 人们创造了新的积分理论——勒贝格 (Lebesgue) 积分理论. 关于这类积分, 我们将在实变函数论这门课程中加以仔细讨论。

要解决满足什么样条件的函数是可积的这个基本的问题, 我们必须从定积分的定义着手. 在定积分的定义中, 分割是任意的, 对于给定的分割, 在每个小区间上的取值也是任意的. 因此, 在函数的可积性的研究中, 我们自然有以下两个问题: 在定积分的定义中,

(a) 是否可取特殊的分割来代替任意的分割?

(b) 能否通过特殊的取值来作和式?

通过本节的讨论, 我们将回答这两个问题。

\subsection{可积的必要条件}

首先我们观察到, 如果一个函数 $f(x)$ 在区间 $[a,b]$ 上无界, 则对于任意的分割
\[
  \Delta:a=x_0<x_1<\cdots<x_n=b,
\]
必存在 $1\leq i_0\leq n$, 使得 $f(x)$ 在 $[x_{i_0-1},x_{i_0}]$ 上无界. 对每个 $i(1\leq i\leq n$ 且 $i\ne i_0)$, 在 $[x_{i-1},x_i]$ 上取定 $\xi_i$, 作和式
\[
  \sum_{i=1,i\ne i_0}^n f(\xi_i)\Delta x_i,
\]
则上述和式在 $\xi_i$ 取定后是一个定数. 记
\[
  M_1=\sum_{i=1,i\ne i_0}^n f(\xi_i)\Delta x_i.
\]
对任意的 $M>0$, 由于 $f(x)$ 在 $[x_{i_0-1},x_{i_0}]$ 上无界, 因此在 $[x_{i_0-1},x_{i_0}]$ 上可取 $\xi_{i_0}$, 使得
\[
  |f(\xi_{i_0})|>\frac{M+|M_1|}{\Delta x_{i_0}},
\]
从而有
\[
\begin{aligned}
  \left|\sum_{i=1}^n f(\xi_i)\Delta x_i\right|
  &=\left|\sum_{i=1,i\ne i_0}^n f(\xi_i)\Delta x_i+f(\xi_{i_0})\Delta x_{i_0}\right|\\
  &\geq |f(\xi_{i_0})|\Delta x_{i_0}-|M_1|\\
  &>M.
\end{aligned}
\]
由于 $M$ 是任意给定的正数, 因此当 $\lambda(\Delta)\to0$ 时, $\sum\limits_{i=1}^n f(\xi_i)\Delta x_i$ 不可能有极限. 所以我们证明了下面的定理。

\noindent\textbf{定理 7.2.1}\quad 若函数 $f(x)\in R[a,b]$, 则 $f(x)$ 在区间 $[a,b]$ 上有界。

在下面的可积性讨论中, 我们总是假定函数 $f(x)$ 是区间 $[a,b]$ 上的有界函数. 一个自然的问题是: 有界函数是否一定可积呢? 此答案是否定的, 如下例所示。

\noindent\textbf{例 7.2.1}\quad 证明狄利克雷函数
\[
  y=D(x)=
  \begin{cases}
    1, & x\in\mathbb Q,\\
    0, & x\in\mathbb R\setminus\mathbb Q
  \end{cases}
\]
在任何区间 $[a,b]$ 上不可积。

\noindent\textbf{证明}\quad 对区间 $[a,b]$ 的任意分割 $\Delta:a=x_0<x_1<\cdots<x_n=b$, 如果在 $[x_{i-1},x_i](i=1,2,\cdots,n)$ 上取一点 $\xi_i(i=1,2,\cdots,n)$, 且 $\xi_i$ 为有理数, 则
\[
  \sum_{i=1}^n D(\xi_i)\Delta x_i=1;
\]
若取 $\xi_i$ 为无理数, 则
\[
  \sum_{i=1}^n D(\xi_i)\Delta x_i=0.
\]
因此, 当 $\lambda(\Delta)\to0$ 时, $\sum\limits_{i=1}^n D(\xi_i)\Delta x_i$ 不可能存在极限. 这说明 $D(x)\notin R[a,b]$。

例 7.2.1 实际上回答了我们在本节开始时提出的关于可积性的第二个问题(问题 (b)), 即在定积分的定义中, 我们不能通过取特殊值的黎曼和来确定函数的可积性。

\subsection{达布理论}

关于可积性的讨论有点类似于序列或函数的上、下极限的讨论. 对区间 $[a,b]$ 上的有界函数 $f(x)$ 和 $[a,b]$ 的一个分割 $\Delta:a=x_0<x_1<\cdots<x_n=b$, 由于 $f(x)$ 在 $[x_{i-1},x_i](i=1,2,\cdots,n)$ 上取值的任意性, 使得其黎曼和具有一定的任意性. 我们注意到: 如果 $f(x)$ 在 $[a,b]$ 上可积, 由定积分的定义可知, 在小区间上任意取值的黎曼和当 $\lambda(\Delta)\to0$ 时有极限, 那么, 对 $\lambda(\Delta)$ 很小时的每个分割 $\Delta$, 所有可能的黎曼和的上确界和下确界必相差不大. 显然 $f(x)$ 在每个小区间上的上确界和下确界只与分割 $\Delta$ 有关, 而与 $\xi_i$ 的取法无关. 为此, 对于分割 $\Delta$ 及 $i=1,2,\cdots,n$, 令
\[
  M_i=\sup_{x\in[x_{i-1},x_i]}\{f(x)\},\qquad
  m_i=\inf_{x\in[x_{i-1},x_i]}\{f(x)\}.
\]
我们作如下的和式:
\[
  \overline S(\Delta)=\sum_{i=1}^n M_i\Delta x_i,\qquad
  \underline S(\Delta)=\sum_{i=1}^n m_i\Delta x_i.
\]
$\overline S(\Delta),\underline S(\Delta)$ 分别称为 $f(x)$ 关于分割 $\Delta$ 的达布 (Darboux) 大和与达布小和。

由于达布大和与达布小和仅仅依赖分割的选取, 对达布大和与达布小和的研究可以避免因 $f(x)$ 在 $[x_{i-1},x_i](i=1,2,\cdots,n)$ 上取值的任意性而使得其黎曼和产生任意性. 如果记 $S(\Delta)=\sum\limits_{i=1}^n f(\xi_i)\Delta x_i$ 为 $f(x)$ 关于分割 $\Delta$ 的任意一个黎曼和, 则显然有
\[
  \underline S(\Delta)\leq S(\Delta)\leq \overline S(\Delta).
\]
另外不难看出, 当 $f(x)\in C[a,b]$ 时, $\underline S(\Delta),\overline S(\Delta)$ 均是黎曼和。

下面我们来研究达布大和与达布小和的一些性质. 若 $\Delta'$ 和 $\Delta''$ 都是区间 $[a,b]$ 的分割, 并且 $\Delta'$ 中的分点均是 $\Delta''$ 中的分点, 我们则称 $\Delta''$ 是 $\Delta'$ 的细分, 记为 $\Delta'\subset\Delta''$。

以下我们总假定函数 $f(x)$ 在区间 $[a,b]$ 上有界, 涉及的分割均是 $[a,b]$ 的分割, 并且所有的黎曼和都是关于 $f(x)$ 的. 对一个给定的分割 $\Delta$, 我们记
\[
  \omega_i=M_i-m_i(i=1,2,\cdots,n).
\]
$\omega_i$ 通常称为 $f(x)$ 在 $[x_{i-1},x_i](i=1,2,\cdots,n)$ 上的振幅. 不难证明
\[
  \omega_i=\sup_{x',x''\in[x_{i-1},x_i]}\{|f(x')-f(x'')|\}.
\]

达布大和与达布小和具有以下性质。

\noindent\textbf{性质 7.2.2}\quad 设 $\Delta'$ 与 $\Delta''$ 为区间 $[a,b]$ 的两个分割, 若 $\Delta'\subset\Delta''$, 则有 $\overline S(\Delta')\geq\overline S(\Delta'')$ 和 $\underline S(\Delta')\leq\underline S(\Delta'')$, 即对于分割的细分, 达布大和不增, 达布小和不减。

\noindent\textbf{证明}\quad 设
\[
  \Delta':a=x_0<x_1<\cdots<x_n=b.
\]
我们先假定
\[
  \Delta'':a=x_0<x_1<\cdots<x_{i-1}<x'<x_i<\cdots<x_n=b,
\]
即 $\Delta''$ 只是在 $\Delta'$ 中加一个分点 $x'$ 所成. 记
\[
  M_i=\sup_{x\in[x_{i-1},x_i]}\{f(x)\},\quad
  M_i'=\sup_{x\in[x_{i-1},x']}\{f(x)\},\quad
  M_i''=\sup_{x\in[x',x_i]}\{f(x)\},
\]
则有 $M_i'\leq M_i$ 和 $M_i''\leq M_i$, 因此有
\[
\begin{aligned}
  \overline S(\Delta')-\overline S(\Delta'')
  &=M_i(x_i-x_{i-1})-[M_i'(x'-x_{i-1})+M_i''(x_i-x')]\\
  &\geq0.
\end{aligned}
\]
对于 $\Delta'$ 中加入多个分点的情形可以比较逐个每次加入一个分点的分割的达布大和并利用上述证明的结论得到. 对于达布小和的情形可类似讨论. 证毕。

\noindent\textbf{性质 7.2.3}\quad 设 $\Delta'$ 和 $\Delta''$ 是区间 $[a,b]$ 的任意两个分割, 则有 $\underline S(\Delta')\leq\overline S(\Delta'')$, 即任意分割的达布小和不大于任意分割的达布大和。

\noindent\textbf{证明}\quad 将 $\Delta',\Delta''$ 中的分点合并在一起构成区间 $[a,b]$ 的一个新的分割, 记其为 $\Delta=\Delta'\cup\Delta''$. 由性质 7.2.2 得
\[
  \underline S(\Delta')\leq \underline S(\Delta)\leq \overline S(\Delta)\leq \overline S(\Delta'').
\]
证毕。

现记
\[
  m=\inf_{x\in[a,b]}\{f(x)\},\qquad M=\sup_{x\in[a,b]}\{f(x)\},
\]
则对任意的分割 $\Delta$, 有
\[
  m(b-a)\leq \underline S(\Delta)\leq \overline S(\Delta)\leq M(b-a).
\]
当分割越来越细时, 性质 7.2.2 告诉我们达布大和与达布小和在某种意义下具有单调性. 因此类似于单调序列极限的性质, 定义
\[
  \underline{\int_a^b} f(x)\mathrm dx=\sup_\Delta\{\underline S(\Delta)\}
  \quad\text{和}\quad
  \overline{\int_a^b} f(x)\mathrm dx=\inf_\Delta\{\overline S(\Delta)\}.
\]
我们分别称 $\underline{\int_a^b} f(x)\mathrm dx$ 和 $\overline{\int_a^b} f(x)\mathrm dx$ 为函数 $f(x)$ 在区间 $[a,b]$ 上的下积分和上积分。

\noindent\textbf{定理 7.2.4(达布定理)}\quad 设函数 $f(x)$ 在区间 $[a,b]$ 上有界, 则
\[
  \lim_{\lambda(\Delta)\to0}\overline S(\Delta)=\overline{\int_a^b} f(x)\mathrm dx,\qquad
  \lim_{\lambda(\Delta)\to0}\underline S(\Delta)=\underline{\int_a^b} f(x)\mathrm dx.
\]

\noindent\textbf{证明}\quad 只证上积分情形, 关于下积分的证明留给读者。

对于 $\forall\varepsilon>0$, 由上积分的定义, 存在一个分割 $\Delta_1:a=x_0'<x_1'<\cdots<x_N'=b(N>1)$, 使得下式成立:
\[
  0\leq \overline S(\Delta_1)-\overline{\int_a^b} f(x)\mathrm dx<\frac{\varepsilon}{2}.
\]
对任意分割 $\Delta:a=x_0<x_1<\cdots<x_n=b$, 由上积分的定义有
\[
  0\leq \overline S(\Delta)-\overline{\int_a^b} f(x)\mathrm dx.
\]
我们记 $\Delta^*=\Delta\cup\Delta_1$, 并考查分割 $\Delta$ 中的每一个小区间 $[x_{i-1},x_i](i=1,2,\cdots,n)$. 如果 $(x_{i-1},x_i)$ 中无 $\Delta_1$ 中的分点, 则在 $\overline S(\Delta)$ 与 $\overline S(\Delta^*)$ 中的相应项同为 $M_i\Delta x_i$. 现在设 $(x_{i-1},x_i)$ 中含有 $\Delta_1$ 中的分点. 记 $\delta_1$ 为 $\Delta_1$ 中 $N$ 个小区间的最小长度, 若预先取 $\lambda(\Delta)<\delta_1$, 则在 $(x_{i-1},x_i)$ 中只有 $\Delta_1$ 中的一个分点 $x'$. 因此 $\overline S(\Delta)$ 与 $\overline S(\Delta^*)$ 中相应项之差为
\[
  M_i(x_i-x_{i-1})-[M_i'(x'-x_{i-1})+M_i''(x_i-x')]
  \leq (M-m)(x_i-x_{i-1})\leq (M-m)\lambda(\Delta),
\]
其中
\[
  M_i=\sup_{x\in[x_{i-1},x_i]}\{f(x)\},
\]
\[
  M_i'=\sup_{x\in[x_{i-1},x']}\{f(x)\},\qquad
  M_i''=\sup_{x\in[x',x_i]}\{f(x)\};
\]
\[
  m=\inf_{x\in[a,b]}\{f(x)\},\qquad M=\sup_{x\in[a,b]}\{f(x)\}.
\]
因此我们得到以下的估计:
\[
  0\leq \overline S(\Delta)-\overline S(\Delta^*)\leq (N-1)(M-m)\lambda(\Delta).
\]
我们不妨设 $m<M$, 否则 $f(x)$ 为常数函数, 此时结论显然成立. 取
\[
  \delta=\min\left\{\delta_1,\frac{\varepsilon}{2(N-1)(M-m)}\right\},
\]
当 $\lambda(\Delta)<\delta$ 时, 有
\[
\begin{aligned}
  0\leq \overline S(\Delta)-\overline{\int_a^b} f(x)\mathrm dx
  &=[\overline S(\Delta)-\overline S(\Delta^*)]+[\overline S(\Delta^*)-\overline S(\Delta_1)]\\
  &\quad+\left[\overline S(\Delta_1)-\overline{\int_a^b} f(x)\mathrm dx\right]\\
  &<\frac{\varepsilon}{2}+0+\frac{\varepsilon}{2}=\varepsilon.
\end{aligned}
\]
证毕。

\noindent\textbf{注}\quad 如果用 ``$\varepsilon$-$\delta$'' 语言来叙述定理 7.2.4, 则它的形式是: 对于 $\forall\varepsilon>0$, $\exists\delta>0$, 对区间 $[a,b]$ 的任意分割 $\Delta$, 当 $\lambda(\Delta)<\delta$ 时, 就有
\[
  0\leq \overline S(\Delta)-\overline{\int_a^b} f(x)\mathrm dx<\varepsilon,\qquad
  0\leq \underline{\int_a^b} f(x)\mathrm dx-\underline S(\Delta)<\varepsilon.
\]

在研究什么样的函数可积这一问题之前, 我们先来考查一下定理 7.2.4 的证明过程. 从该定理的证明中读者不难发现以下事实: 对于 $\forall\varepsilon>0$, 若我们能找到一个特别的分割 $\Delta_1$, 使得
\[
  0\leq \overline S(\Delta_1)-\overline{\int_a^b} f(x)\mathrm dx<\frac{\varepsilon}{2},
\]
则必定存在 $\delta>0$, 使得对任意分割 $\Delta$, 当 $\lambda(\Delta)<\delta$ 时, 有
\[
  0\leq \overline S(\Delta)-\overline{\int_a^b} f(x)\mathrm dx<\varepsilon.
\]
上述事实告诉我们如下结论: $\lim\limits_{\lambda(\Delta)\to0}\overline S(\Delta)=\overline{\int_a^b} f(x)\mathrm dx$ 的充分必要条件是: 对于 $\forall\varepsilon>0$, 存在分割 $\Delta$, 使得
\[
  0\leq \overline S(\Delta)-\overline{\int_a^b} f(x)\mathrm dx<\varepsilon.
\]

同理, 对于下积分, 我们也有相应的性质. 上积分与下积分的这个性质将应用到可积性问题的讨论中。

现在我们来证明下面的定理。

\noindent\textbf{定理 7.2.5}\quad 设函数 $f(x)$ 在区间 $[a,b]$ 上有界, 则 $f(x)\in R[a,b]$ 的充分必要条件是
\[
  \underline{\int_a^b} f(x)\mathrm dx=\overline{\int_a^b} f(x)\mathrm dx.
\]

\noindent\textbf{证明}\quad 必要性\quad 我们观察到, 对任意的分割, 通过在每个小区间上的特殊取值, 可以使黎曼和任意地接近达布大和或达布小和. 设 $f(x)\in R[a,b]$, 则对于 $\forall\varepsilon>0$, $\exists\delta>0$, 对 $[a,b]$ 的任意分割
\[
  \Delta:a=x_0<x_1<x_2<\cdots<x_n=b,
\]
只要 $\lambda(\Delta)<\delta$, 就有
\[
  \left|\sum_{i=1}^n f(\xi_i)\Delta x_i-\int_a^b f(x)\mathrm dx\right|<\varepsilon,
\]
其中 $\xi_i\in[x_{i-1},x_i](i=1,2,\cdots,n)$。
由上、下确界的定义, 存在 $\xi_i',\xi_i''\in[x_{i-1},x_i](i=1,2,\cdots,n)$, 使得
\[
  f(\xi_i')\geq M_i-\frac{\varepsilon}{b-a}\quad\text{及}\quad f(\xi_i'')\leq m_i+\frac{\varepsilon}{b-a},
\]
其中
\[
  M_i=\sup_{x\in[x_{i-1},x_i]}\{f(x)\},\qquad
  m_i=\inf_{x\in[x_{i-1},x_i]}\{f(x)\}\quad(i=1,2,\cdots,n).
\]
因此有
\[
  \overline{\int_a^b} f(x)\mathrm dx-\varepsilon
  \leq \sum_{i=1}^n M_i\Delta x_i-\varepsilon
  \leq \sum_{i=1}^n f(\xi_i')\Delta x_i
  < \int_a^b f(x)\mathrm dx+\varepsilon
\]
和
\[
  \underline{\int_a^b} f(x)\mathrm dx+\varepsilon
  \geq \sum_{i=1}^n m_i\Delta x_i+\varepsilon
  \geq \sum_{i=1}^n f(\xi_i'')\Delta x_i
  > \int_a^b f(x)\mathrm dx-\varepsilon.
\]
由此得
\[
  0\leq \overline{\int_a^b} f(x)\mathrm dx-\underline{\int_a^b} f(x)\mathrm dx<4\varepsilon.
\]
由于 $\overline{\int_a^b} f(x)\mathrm dx-\underline{\int_a^b} f(x)\mathrm dx$ 是固定常数, 由 $\varepsilon$ 的任意性, 必有
\[
  \underline{\int_a^b} f(x)\mathrm dx=\overline{\int_a^b} f(x)\mathrm dx.
\]
必要性得证。

\noindent\textbf{充分性}\quad 由定理 7.2.4 可知, 对于 $\forall\varepsilon>0$, $\exists\delta>0$, 对于任意的分割 $\Delta$, 当 $\lambda(\Delta)<\delta$ 时, 对于 $\Delta$ 的任意黎曼和 $S(\Delta)$, 有
\[
  \underline{\int_a^b} f(x)\mathrm dx-\varepsilon
  \leq \underline S(\Delta)\leq S(\Delta)\leq \overline S(\Delta)
  \leq \overline{\int_a^b} f(x)\mathrm dx+\varepsilon.
\]
记 $I=\underline{\int_a^b} f(x)\mathrm dx=\overline{\int_a^b} f(x)\mathrm dx$, 则由上式得
\[
  |S(\Delta)-I|<\varepsilon.
\]
由定积分的定义知 $f(x)\in R[a,b]$. 充分性得证. 证毕。

为了叙述简便, 下面对于 $[a,b]$ 的一个分割 $\Delta:a=x_0<x_1<\cdots<x_n=b$, 我们总记 $\omega_i$ 为 $f(x)$ 在区间 $[x_{i-1},x_i]$ 上的振幅。
由定理 7.2.4 和后面的注及定理 7.2.5, 我们有下面的结果。

\noindent\textbf{定理 7.2.6}\quad 设函数 $f(x)$ 在区间 $[a,b]$ 上有界, 则下面的三个结论相互等价:

(1) $f(x)$ 在区间 $[a,b]$ 上可积;

(2) 对于 $\forall\varepsilon>0$, 存在区间 $[a,b]$ 的一个分割 $\Delta$, 使得
\[
  \sum_{i=1}^n \omega_i\Delta x_i<\varepsilon;
\]

(3) 对于 $\forall\varepsilon>0$, $\forall\sigma>0$, 存在区间 $[a,b]$ 的一个分割 $\Delta$, 使得 $\omega_i\geq\varepsilon$ 的小区间 $[x_{i-1},x_i]$ 的长度总和小于 $\sigma$。

\noindent\textbf{证明}\quad 我们先证 (1)$\Rightarrow$(2). 设 $f(x)\in R[a,b]$, 由定理 7.2.5 我们有
\[
  \underline{\int_a^b} f(x)\mathrm dx=\overline{\int_a^b} f(x)\mathrm dx.
\]
再由上、下积分的定义易知, 对于 $\forall\varepsilon>0$, 存在 $[a,b]$ 的分割 $\Delta$, 使得
\[
  \underline{\int_a^b} f(x)\mathrm dx-\frac{\varepsilon}{2}<\underline S(\Delta)\leq \overline S(\Delta)<\overline{\int_a^b} f(x)\mathrm dx+\frac{\varepsilon}{2}.
\]
因此有
\[
  \sum_{i=1}^n \omega_i\Delta x_i=\overline S(\Delta)-\underline S(\Delta)<\varepsilon.
\]
(2) 得证。

现证 (2)$\Rightarrow$(1). 由于对于 $\forall\varepsilon>0$, 存在 $[a,b]$ 的一个分割 $\Delta:a=x_0<x_1<\cdots<x_n=b$, 使得
\[
  \sum_{i=1}^n \omega_i\Delta x_i<\varepsilon,
\]
因此有
\[
  0\leq \overline{\int_a^b} f(x)\mathrm dx-\underline{\int_a^b} f(x)\mathrm dx
  \leq \overline S(\Delta)-\underline S(\Delta)
  =\sum_{i=1}^n\omega_i\Delta x_i<\varepsilon.
\]
由 $\varepsilon$ 的任意性, 我们有
\[
  \underline{\int_a^b} f(x)\mathrm dx=\overline{\int_a^b} f(x)\mathrm dx.
\]
再由定理 7.2.5 我们知 $f(x)\in R[a,b]$。

下面我们来证明 (2)$\Longleftrightarrow$(3). 首先注意到, 对于 $\forall\varepsilon>0$, 对 $[a,b]$ 的分割
\[
  \Delta:a=x_0<x_1<\cdots<x_n=b,
\]
总有
\[
  \sum_{i=1}^n\omega_i\Delta x_i=\sum_{(1)}\omega_i\Delta x_i+\sum_{(2)}\omega_i\Delta x_i,
\]
其中 $\sum_{(1)}$ 表示对 $\omega_i<\varepsilon$ 的求和, $\sum_{(2)}$ 表示对 $\omega_i\geq\varepsilon$ 的求和。
先证 (2)$\Rightarrow$(3). 对于 $\forall\varepsilon>0$, $\forall\sigma>0$, 由 (2), 存在分割
\[
  \Delta:a=x_0<x_1<\cdots<x_n=b,
\]
使得
\[
  \sum_{i=1}^n\omega_i\Delta x_i<\sigma\varepsilon.
\]
因此有
\[
  \varepsilon\sum_{(2)}\Delta x_i\leq \sum_{(2)}\omega_i\Delta x_i\leq \sum_{i=1}^n\omega_i\Delta x_i<\sigma\varepsilon,
\]
即
\[
  \sum_{(2)}\Delta x_i<\sigma.
\]

再证 (3)$\Rightarrow$(2). 记 $\omega$ 为 $f(x)$ 在 $[a,b]$ 上的振幅, 不妨设 $f(x)$ 不是常数函数, 因此 $\omega>0$. 由 (3), 对于 $\dfrac{\varepsilon}{2(b-a)}$ 及 $\sigma=\dfrac{\varepsilon}{2\omega}$, 存在分割
\[
  \Delta:a=x_0<x_1<\cdots<x_n=b,
\]
使得 $\omega_i\geq\dfrac{\varepsilon}{2(b-a)}$ 的小区间 $[x_{i-1},x_i]$ 长度的总和小于 $\sigma$. 现在记 $\sum_{(1)'}\omega_i\Delta x_i$ 是对 $\omega_i<\dfrac{\varepsilon}{2(b-a)}$ 的求和, $\sum_{(2)'}\omega_i\Delta x_i$ 是对 $\omega_i\geq\dfrac{\varepsilon}{2(b-a)}$ 的求和, 则有
\[
\begin{aligned}
  \sum_{i=1}^n\omega_i\Delta x_i
  &=\sum_{(1)'}\omega_i\Delta x_i+\sum_{(2)'}\omega_i\Delta x_i\\
  &<\frac{\varepsilon}{2(b-a)}\sum_{(1)'}\Delta x_i+\omega\sum_{(2)'}\Delta x_i\\
  &<\frac{\varepsilon}{2(b-a)}(b-a)+\omega\frac{\varepsilon}{2\omega}
  =\varepsilon.
\end{aligned}
\]
证毕。

\noindent\textbf{注}\quad 定理 7.2.6 实际上回答了我们在本节开始时提出的第一个关于可积性的问题(问题 (a)), 即在定积分的定义中, 我们可以取一列特殊的分割 $\Delta_n$, 只要当 $\lambda(\Delta_n)\to0$ 时, $f(x)$ 关于 $\Delta_n$ 的任意黎曼和可以任意接近一个常数 $I$, 则 $f(x)\in R[a,b]$ 且 $I=\int_a^b f(x)\mathrm dx$。

\subsection{可积函数类}

由上一节的结果, 我们可以找到一些常见的可积函数类。

\noindent\textbf{定理 7.2.7}\quad 设函数 $f(x)$ 在区间 $[a,b]$ 上有界, 且只有有限个间断点, 则 $f(x)\in R[a,b]$。

\noindent\textbf{证明}\quad 记 $f(x)$ 在 $[a,b]$ 上的振幅为 $\omega$, 我们不妨设 $\omega>0$(否则 $f(x)$ 为一个常数函数). 再设 $f(x)$ 在 $[a,b]$ 上的间断点的个数为 $N$. 对于 $\forall\varepsilon>0$, 在 $[a,b]$ 上取 $N$ 个互不相交的闭区间, 使得满足: (1) 每个区间的长度大于零小于 $\dfrac{\varepsilon}{2N\omega}$; (2) 它们包含了 $f(x)$ 在 $[a,b]$ 上的所有的间断点; (3) 它们的内部(设 $[c,d]$ 为一闭区间, 称 $(c,d)$ 为该区间的内部)包含了 $f(x)$ 在 $(a,b)$ 内所有的间断点. $[a,b]$ 除去这组区间后得到的每个小区间加上端点则为有限个闭区间: $I_1,I_2,\cdots,I_K$. 由于 $f(x)$ 在每个小闭区间 $I_j(j=1,2,\cdots,K)$ 上连续, 从而一致连续, 因此存在 $\delta>0$, 使得当 $x',x''$ 属于同一个闭区间 $I_j$, 且 $|x'-x''|<\delta$ 时, 有
\[
  |f(x')-f(x'')|<\frac{\varepsilon}{2(b-a)}.
\]
将每个小区间 $I_j$ 适当等分, 使得等分后的小区间长度小于 $\delta$, 则所有这些小区间的等分点与原来的 $N$ 个区间的端点构成 $[a,b]$ 的一个分割 $\Delta:a=x_0<x_1<\cdots<x_n=b$. 对此分割有
\[
  \sum_{i=1}^n\omega_i\Delta x_i< \frac{\varepsilon}{2N\omega}\omega N+\frac{\varepsilon(b-a)}{2(b-a)}=\varepsilon.
\]
证毕。

另外, 由定积分的定义以及以上我们得到的可积性的判别法容易看出, 若函数 $f(x)\in R[a,b]$, 函数 $g(x)$ 在 $[a,b]$ 上有定义且与 $f(x)$ 在 $[a,b]$ 上只有有限多个点上的取值不同, 则 $g(x)\in R[a,b]$, 并且有
\[
  \int_a^b f(x)\mathrm dx=\int_a^b g(x)\mathrm dx.
\]

\noindent\textbf{定理 7.2.8}\quad 设函数 $f(x)$ 在区间 $[a,b]$ 上单调, 则 $f(x)\in R[a,b]$。

\noindent\textbf{证明}\quad 不妨设 $f(x)$ 在 $[a,b]$ 上单调上升. 对于 $\forall\varepsilon>0$, 将 $[a,b]$ 等分成 $n$ 个小区间, 使得
\[
  \frac{(b-a)[f(b)-f(a)]}{n}<\varepsilon.
\]
设其对应的分割为 $a=x_0<x_1<\cdots<x_n=b$, 则有
\[
\begin{aligned}
  \sum_{i=1}^n \omega_i\Delta x_i
  &=\sum_{i=1}^n [f(x_i)-f(x_{i-1})]\frac{b-a}{n}\\
  &=\frac{b-a}{n}[f(b)-f(a)]<\varepsilon.
\end{aligned}
\]
证毕。

\noindent\textbf{例 7.2.2}\quad 证明黎曼函数
\[
R(x)=
\begin{cases}
  \dfrac{1}{p}, & x\in[0,1],\ x=\dfrac{q}{p}(p,q>0,\text{为既约正整数}),\\
  0, & x\in(0,1)\setminus\mathbb Q,\\
  1, & x=0
\end{cases}
\]
在区间 $[0,1]$ 上可积, 且 $\int_0^1 R(x)\mathrm dx=0$。

\noindent\textbf{证明}\quad 对于 $\forall\varepsilon>0$, $\forall\sigma>0$, 在 $[0,1]$ 中仅存在有限多个点使得 $R(x)\geq\varepsilon$. 设这有限多个点为 $x_1',\cdots,x_N'$. 在 $[0,1]$ 中取 $N$ 个互不相交的小闭区间, 使得每一个区间的长度大于零小于 $\dfrac{\sigma}{N}$, 并且它们包含这 $N$ 个点. 进一步我们要求当 $x_i'\neq0,1$ 时, $x_i'$ 落在某个小区间的内部. 对于这 $N$ 个小区间的端点构成 $[0,1]$ 的分割 $\Delta$, 那些振幅大于或等于 $\varepsilon$ 的小区间的总和显然小于 $\dfrac{\sigma}{N}\cdot N=\sigma$. 由定理 7.2.6 的 (3) 知 $R(x)\in R[0,1]$。

由于 $R(x)$ 在 $[0,1]$ 上可积, 并且对任意的分割
\[
  \Delta:a=x_0<x_1<\cdots<x_n=b,
\]
在每个 $[x_{i-1},x_i]$ 中取 $\xi_i$ 为无理数, 则其黎曼和为
\[
  \sum_{i=1}^n R(\xi_i)\Delta x_i=0,
\]
因此
\[
  \int_0^1 R(x)\mathrm dx=0.
\]

\noindent\textbf{例 7.2.3}\quad 证明函数
\[
f(x)=
\begin{cases}
  \sin\dfrac{1}{x}, & x\neq0,\\
  0, & x=0
\end{cases}
\]
在任何区间 $[a,b]$ 上可积。

\noindent\textbf{证明}\quad 当 $0\notin[a,b]$ 时, $f(x)$ 在 $[a,b]$ 上连续, 从而 $f(x)\in R[a,b]$; 当 $0\in[a,b]$ 时, $f(x)$ 在 $[a,b]$ 上有界, 且只有一个间断点 $x_0=0$, 因此由定理 7.2.7 知 $f(x)\in R[a,b]$。

从本例中可以看出, 尽管函数 $f(x)=\sin\dfrac{1}{x}$ 在 $x_0=0$ 的邻域中变化非常剧烈, 但它仍然在包含 $x_0$ 的区间上可积。

函数的可积性问题将在实变函数论课程中得到完全解决. 在实变函数论课程中将证明以下结论。

\noindent\textbf{定理 7.2.9(勒贝格定理)}\quad 设函数 $f(x)$ 在区间 $[a,b]$ 上有界, 记 $E$ 为 $f(x)$ 在 $[a,b]$ 上的间断点集, 则 $f(x)\in R[a,b]$ 的充分必要条件是: 对于 $\forall\varepsilon>0$, 存在一列开区间 $(x_{i-1},x_i)(i=1,2,\cdots)$, 使得
\[
  E\subset \bigcup_{i=1}^{\infty}(x_{i-1},x_i),
  \qquad
  \text{并且对于 } \forall n\in\mathbb N,\ \text{有 } \sum_{i=1}^n(x_i-x_{i-1})<\varepsilon.
\]

\section{定积分的性质}

前面我们已经对可积性问题进行了讨论, 在本节中将讨论定积分的一些基本性质。

首先, 我们作两条自然的规定:
\[
  \text{(1) }\int_a^a f(x)\mathrm dx=0;
\]
\[
  \text{(2) 当函数 } f(x) \text{ 在区间 } [a,b] \text{ 上可积时, }
  \int_b^a f(x)\mathrm dx=-\int_a^b f(x)\mathrm dx.
\]
定积分具有以下基本性质。

\noindent\textbf{定理 7.3.1(线性性质)}\quad 设函数 $f(x),g(x)\in R[a,b]$, $\alpha,\beta\in\mathbb R$, 则 $\alpha f(x)+\beta g(x)\in R[a,b]$, 并且有
\[
  \int_a^b[\alpha f(x)+\beta g(x)]\mathrm dx
  =\alpha\int_a^b f(x)\mathrm dx+\beta\int_a^b g(x)\mathrm dx.
\]
此性质可由定积分的定义直接推出, 请读者自证。

\noindent\textbf{定理 7.3.2}\quad 设函数 $f(x)\in R[a,b]$, 则 $|f(x)|\in R[a,b]$, 并且有
\[
  \left|\int_a^b f(x)\mathrm dx\right|\leq \int_a^b |f(x)|\mathrm dx.
\]

\noindent\textbf{证明}\quad 对于 $\forall\varepsilon>0$, 由 $f(x)\in R[a,b]$, 存在 $[a,b]$ 的分割
\[
  \Delta:a=x_0<x_1<\cdots<x_n=b,
\]
记 $\omega_i(f)$ 为 $f(x)$ 在 $[x_{i-1},x_i](i=1,2,\cdots,n)$ 上的振幅, 则有
\[
  \sum_{i=1}^n \omega_i(f)\Delta x_i<\varepsilon.
\]
现记 $\omega_i(|f|)$ 为 $|f(x)|$ 在 $[x_{i-1},x_i](i=1,2,\cdots,n)$ 上的振幅, 由于当 $x',x''\in[x_{i-1},x_i]$ 时, 有
\[
  \bigl||f(x')|-|f(x'')|\bigr|\leq |f(x')-f(x'')|,
\]
因此有
\[
  \omega_i(|f|)\leq \omega_i(f)\quad(i=1,2,\cdots,n),
\]
从而有
\[
  \sum_{i=1}^n \omega_i(|f|)\Delta x_i<\varepsilon.
\]
所以 $|f(x)|\in R[a,b]$。

对 $[a,b]$ 的任意分割 $\Delta$ 及 $[x_{i-1},x_i](i=1,2,\cdots,n)$ 上任取的 $\xi_i$, 我们有
\[
  \left|\sum_{i=1}^n f(\xi_i)\Delta x_i\right|
  \leq \sum_{i=1}^n |f(\xi_i)|\Delta x_i.
\]
因此, 当 $\lambda(\Delta)\to0$ 时, 有
\[
  \left|\int_a^b f(x)\mathrm dx\right|\leq \int_a^b |f(x)|\mathrm dx.
\]
证毕。

\noindent\textbf{定理 7.3.3}\quad 设 $a<c<b$, 则函数 $f(x)\in R[a,b]$ 的充分必要条件是: $f(x)\in R[a,c]$ 和 $f(x)\in R[c,b]$. 当 $f(x)\in R[a,b]$ 时, 下式成立:
\begin{equation}
  \int_a^b f(x)\mathrm dx=\int_a^c f(x)\mathrm dx+\int_c^b f(x)\mathrm dx.
  \tag{7.3.1}
\end{equation}

\noindent\textbf{证明}\quad 必要性\quad 设函数 $f(x)\in R[a,b]$, 我们证明 $f(x)$ 在 $[a,b]$ 的任何子区间 $[a_1,b_1]\subset[a,b]$ 上可积. 事实上, 由于 $f(x)$ 在 $[a,b]$ 上可积, 对于 $\forall\varepsilon>0$, 存在 $[a,b]$ 的分割
\[
  \Delta:a=x_0<x_1<\cdots<x_n=b,
\]
记 $\omega_i$ 为 $f(x)$ 在 $[x_{i-1},x_i](i=1,2,\cdots,n)$ 上的振幅, 则有
\[
  \sum_{i=1}^n \omega_i\Delta x_i<\varepsilon.
\]
由于 $\Delta$ 落在 $[a_1,b_1]$ 中的分割点与 $a_1,b_1$ 构成了 $[a_1,b_1]$ 的一个分割
\[
  \Delta':a_1=x_0'<x_1'<\cdots<x_m'=b_1,
\]
对于此分割显然有
\[
  \sum_{i=1}^m \omega_i'\Delta x_i'<\varepsilon,
\]
其中 $\omega_i'$ 为 $f(x)$ 在 $[x_{i-1}',x_i']$ 上的振幅, 而 $\Delta x_i'=x_i'-x_{i-1}'(i=1,2,\cdots,m)$, 因此 $f(x)$ 在 $[a_1,b_1]$ 上可积. 特别地, 有 $f(x)\in R[a,c]$, $f(x)\in R[c,b]$。

充分性\quad 对于 $\forall\varepsilon>0$, 由 $f(x)\in R[a,c]$, 存在 $[a,c]$ 的分割
\[
  \Delta_1:a=x_0<x_1<\cdots<x_n=c,
\]
记 $\omega_i$ 为 $f(x)$ 在 $[x_{i-1},x_i](i=1,2,\cdots,n)$ 上的振幅, 则有
\[
  \sum_{i=1}^n\omega_i\Delta x_i<\frac{\varepsilon}{2}.
\]
由 $f(x)\in R[c,b]$, 存在 $[c,b]$ 上的分割
\[
  \Delta_2:c=x_0'<x_1'<\cdots<x_m'=b,
\]
记 $\omega_i'$ 为 $f(x)$ 在 $[x_{i-1}',x_i'](i=1,2,\cdots,m)$ 上的振幅及 $\Delta x_i'=x_i'-x_{i-1}'$, 则有
\[
  \sum_{i=1}^m\omega_i'\Delta x_i'<\frac{\varepsilon}{2}.
\]
显然, $\Delta_1\cup\Delta_2$ 是 $[a,b]$ 的一个分割 $\Delta$, 且有
\[
  \sum_{i=1}^n\omega_i\Delta x_i+\sum_{i=1}^m\omega_i'\Delta x_i'<\varepsilon.
\]
而上式左边刚好是 $f(x)$ 关于 $\Delta$ 的每个小区间的振幅乘以对应小区间的长度之和, 因此 $f(x)\in R[a,b]$。

现证定积分等式 (7.3.1) 成立. 事实上, 对 $[a,c]$ 的任意分割
\[
  \Delta_1:a=x_0<x_1<\cdots<x_n=c
\]
作黎曼和 $\sum_{i=1}^n f(\xi_i)\Delta x_i$, 对 $[c,b]$ 的任意分割
\[
  \Delta_2:c=x_0'<x_1'<\cdots<x_m'=b
\]
作黎曼和 $\sum_{i=1}^m f(\xi_i')\Delta x_i'$, 则
\[
  \sum_{i=1}^n f(\xi_i)\Delta x_i+\sum_{i=1}^m f(\xi_i')\Delta x_i'
\]
是 $f(x)$ 在 $[a,b]$ 上关于分割 $\Delta_1\cup\Delta_2$ 的黎曼和. 由于以上的各个黎曼和当小区间最大长度趋于零时都存在极限, 因此定积分等式 (7.3.1) 成立. 证毕。

关于本性质我们要指出的是: 如果不假定 $a<c<b$, 但假定函数 $f(x)$ 在所涉及的区间是可积的, 利用在本节开始时的两个规定, 则等式 (7.3.1) 总是成立的. 例如, 设 $a<b<c$, 假定 $f(x)\in R[a,c]$, 则仍然成立
\[
  \int_a^b f(x)\mathrm dx=\int_a^c f(x)\mathrm dx+\int_c^b f(x)\mathrm dx.
\]
事实上, 我们有
\[
\begin{aligned}
  \int_a^b f(x)\mathrm dx
  &=\int_a^b f(x)\mathrm dx+\int_b^c f(x)\mathrm dx-\int_b^c f(x)\mathrm dx\\
  &=\int_a^c f(x)\mathrm dx+\int_c^b f(x)\mathrm dx.
\end{aligned}
\]

\noindent\textbf{定理 7.3.4}\quad 设函数 $f(x),g(x)\in R[a,b]$, 并且对于 $\forall x\in[a,b]$, 有 $f(x)\geq g(x)$, 则
\[
  \int_a^b f(x)\mathrm dx\geq \int_a^b g(x)\mathrm dx.
\]
定理 7.3.4 的证明是简单的. 事实上对函数 $f(x)-g(x)$ 用定积分的定义及定积分的线性性质, 即可推出该定理的结论。

设函数 $g(x)$ 在区间 $[a,b]$ 上有定义, 如果 $[a,b]$ 是有限个两两互不相交的小区间的并, 并且 $g(x)$ 在每个小区间上均为常数函数, 则称 $g(x)$ 为 $[a,b]$ 上的一个阶梯函数. 显然闭区间 $[a,b]$ 上的任何阶梯函数是有界的, 而且至多只有有限个间断点, 因此阶梯函数是可积的。

\noindent\textbf{例 7.3.1}\quad 设函数 $f(x)\in R[a,b]$, 证明: 对于 $\forall\varepsilon>0$,

(1) 存在 $[a,b]$ 上的阶梯函数 $h(x)$, 使得 $\int_a^b |f(x)-h(x)|\mathrm dx<\varepsilon$;

(2) 存在 $[a,b]$ 上的连续函数 $g(x)$, 使得 $\int_a^b |f(x)-g(x)|\mathrm dx<\varepsilon$。

\noindent\textbf{证明}\quad 由于 $f(x)\in R[a,b]$, 对于 $\forall\varepsilon>0$, 存在 $[a,b]$ 的分割
\[
  \Delta:a=x_0<x_1<\cdots<x_n=b,
\]
记 $M_i,m_i$ 和 $\omega_i$ 分别为 $f(x)$ 在 $[x_{i-1},x_i](i=1,2,\cdots,n)$ 上的上、下确界和振幅, 则有
\[
  \sum_{i=1}^n\omega_i\Delta x_i<\varepsilon.
\]
(1) 定义阶梯函数
\[
  \overline f(x)=
  \begin{cases}
    M_i, & x\in[x_{i-1},x_i),\ i=1,2,\cdots,n-1,\\
    M_n, & x\in[x_{n-1},b]
  \end{cases}
\]
和
\[
  \underline f(x)=
  \begin{cases}
    m_i, & x\in[x_{i-1},x_i),\ i=1,2,\cdots,n-1,\\
    m_n, & x\in[x_{n-1},b],
  \end{cases}
\]
则 $\overline f(x)$ 与 $\underline f(x)$ 均为区间 $[a,b]$ 上的阶梯函数. 为了证明 $\overline f(x)$ 与 $\underline f(x)$ 满足 (1) 中的结论, 我们令
\[
  \overline f_1(x)=
  \begin{cases}
    \overline f(x), & x\in[a,b]\text{ 且 }x\neq x_i,\\
    f(x_i), & x=x_i
  \end{cases}
  \quad(i=0,1,\cdots,n).
\]
\[
  \underline f_1(x)=
  \begin{cases}
    \underline f(x), & x\in[a,b]\text{ 且 }x\neq x_i,\\
    f(x_i), & x=x_i
  \end{cases}
  \quad(i=0,1,\cdots,n).
\]
容易看出 $\overline f(x)$, $\overline f_1(x)$, $\underline f(x)$ 与 $\underline f_1(x)$ 均为 $[a,b]$ 上的可积函数, 且有
\[
  \int_a^b |f(x)-\overline f(x)|\mathrm dx
  =\int_a^b |f(x)-\overline f_1(x)|\mathrm dx
\]
和
\[
  \int_a^b |f(x)-\underline f(x)|\mathrm dx
  =\int_a^b |f(x)-\underline f_1(x)|\mathrm dx.
\]
由振幅的定义, 对于 $\forall x\in[x_{i-1},x_i](i=1,2,\cdots,n)$, 有
\begin{equation}
  |f(x)-\overline f_1(x)|\leq\omega_i
  \tag{7.3.2}
\end{equation}
和
\begin{equation}
  |f(x)-\underline f_1(x)|\leq\omega_i,
  \tag{7.3.3}
\end{equation}
从而有
\[
\begin{aligned}
  \int_a^b |f(x)-\underline f(x)|\mathrm dx
  &=\int_a^b |f(x)-\underline f_1(x)|\mathrm dx\\
  &=\sum_{i=1}^n\int_{x_{i-1}}^{x_i}|f(x)-\underline f_1(x)|\mathrm dx\\
  &\leq \sum_{i=1}^n\omega_i\Delta x_i<\varepsilon
\end{aligned}
\]
和
\[
  \int_a^b |f(x)-\overline f(x)|\mathrm dx<\varepsilon.
\]
取 $h(x)=\underline f(x)$(或 $\overline f(x)$), 即得 (1)。

(2) 作折线依次连接 $(x_{i-1},f(x_{i-1}))$ 和 $(x_i,f(x_i))(i=1,2,\cdots,n)$. 设 $g(x)$ 为 $[a,b]$ 上的函数, 其图像恰为该折线, 则 $g(x)$ 在 $[a,b]$ 上连续($g(x)$ 也称为是分段线性函数)。

同理由振幅的定义, 对于 $\forall x\in[x_{i-1},x_i](i=1,2,\cdots,n)$, 有
\begin{equation}
  |f(x)-g(x)|\leq\omega_i,
  \tag{7.3.4}
\end{equation}
因此有
\[
  \int_a^b |f(x)-g(x)|\mathrm dx<\varepsilon,
\]
从而 (2) 得证。

\noindent\textbf{注}\quad 在第十二章傅里叶(Fourier)级数的研究中, 我们还要用到以下结论: 设函数 $f(x)\in R[a,b]$, 则对于 $\forall\varepsilon>0$, 存在区间 $[a,b]$ 上的连续函数 $g(x)$, 使得
\[
  \int_a^b [f(x)-g(x)]^2\mathrm dx<\varepsilon.
\]
利用例 7.3.1, 我们很容易证明此结论. 事实上, 由 $f(x)\in R[a,b]$, 从而存在 $M>0$, 使得对于 $x\in[a,b]$, 有 $|f(x)|\leq M$. 对于 $\forall\varepsilon>0$, 由例 7.3.1, 我们可以找到 $[a,b]$ 上的连续函数 $g(x)$, 使得
\[
  \int_a^b |f(x)-g(x)|\mathrm dx<\frac{\varepsilon}{2M}.
\]
显然, 从例 7.3.1 中 $g(x)$ 的构造不难看出, 对于 $\forall x\in[a,b]$, 有 $|g(x)|\leq M$, 因此有
\[
\begin{aligned}
  \int_a^b [f(x)-g(x)]^2\mathrm dx
  &=\int_a^b |f(x)-g(x)|^2\mathrm dx\\
  &\leq \int_a^b (|f(x)|+|g(x)|)|f(x)-g(x)|\mathrm dx\\
  &\leq \int_a^b 2M|f(x)-g(x)|\mathrm dx\\
  &<2M\cdot\frac{\varepsilon}{2M}=\varepsilon.
\end{aligned}
\]

我们下面继续讨论定积分的性质。

\noindent\textbf{定理 7.3.5}\quad 设函数 $f(x),g(x)\in R[a,b]$, 则 $f(x)g(x)\in R[a,b]$。

\noindent\textbf{证明}\quad 由 $f(x),g(x)$ 的可积性知, $f(x)$ 与 $g(x)$ 均是有界函数. 我们取 $M>0$, 使得当 $x\in[a,b]$ 时, 有 $|f(x)|\leq M$ 及 $|g(x)|\leq M$. 对于 $\forall\varepsilon>0$, 再由 $f(x),g(x)$ 的可积性, 存在 $[a,b]$ 的分割
\[
  \Delta:a=x_0<x_1<\cdots<x_n=b,
\]
记 $\omega_i(h)$ 为函数 $h(x)$ 在 $[x_{i-1},x_i](i=1,2,\cdots,n)$ 上的振幅, 则有
\begin{equation}
  \sum_{i=1}^n \omega_i(f)\Delta x_i<\frac{\varepsilon}{2M}
  \tag{7.3.5}
\end{equation}
及
\begin{equation}
  \sum_{i=1}^n \omega_i(g)\Delta x_i<\frac{\varepsilon}{2M}.
  \tag{7.3.6}
\end{equation}
对任意的 $x',x''\in[x_{i-1},x_i]$, 由
\[
\begin{aligned}
  |f(x'')g(x'')-f(x')g(x')|
  &\leq |f(x'')g(x'')-f(x'')g(x')|\\
  &\quad+|f(x'')g(x')-f(x')g(x')|\\
  &\leq |f(x'')||g(x'')-g(x')|+|g(x')||f(x'')-f(x')|
\end{aligned}
\]
推出
\[
  \omega_i(fg)\leq M[\omega_i(f)+\omega_i(g)]
  \quad(i=1,2,\cdots,n).
\]
因此
\[
  \sum_{i=1}^n\omega_i(fg)\Delta x_i
  \leq M\left[\sum_{i=1}^n\omega_i(f)\Delta x_i+\sum_{i=1}^n\omega_i(g)\Delta x_i\right]<\varepsilon.
\]
证毕。

容易看出, 当 $f(x),g(x)\in R[a,b]$ 时, $\dfrac{f(x)}{g(x)}$ 未必在 $[a,b]$ 上可积(如此时 $\dfrac{f(x)}{g(x)}$ 可以是无界函数); 若 $f(x)$ 与 $g(x)$ 可以复合, $f(g(x))$ 在 $[a,b]$ 上也不一定可积. 例如, 取 $g(x)=R(x)$, 其中 $R(x)$ 是区间 $[0,1]$ 上的黎曼函数, 而取
\[
  f(x)=
  \begin{cases}
    1, & x\neq0,\\
    0, & x=0,
  \end{cases}
\]
则 $f(g(x))$ 为区间 $[0,1]$ 上的狄利克雷函数. 因此 $f(g(x))$ 在 $[0,1]$ 上不可积. 但当 $f(x)$ 连续时, 若 $g(x)$ 是 $[a,b]$ 上的可积函数并且 $f(g(x))$ 在 $[a,b]$ 上有定义, 则必定有 $f(g(x))\in R[a,b]$(见本章习题)。

\noindent\textbf{例 7.3.2}\quad 设函数 $f(x)$ 在区间 $[a,b]$ 上连续, 证明:

(1) 如果对 $[a,b]$ 的任何子区间 $[a_1,b_1]$ 有 $\int_{a_1}^{b_1}f(x)\mathrm dx=0$, 则 $f(x)\equiv0,\ x\in[a,b]$;

(2) 若 $f(x)\geq0$, 并且 $\int_a^b f(x)\mathrm dx=0$, 则 $f(x)\equiv0,\ x\in[a,b]$。

\noindent\textbf{证明}\quad (1) 用反证法. 倘若 $f(x)$ 不恒为 $0$, 则存在 $x_0\in[a,b]$, 使得 $f(x_0)\neq0$. 不妨设 $f(x_0)>0$, 且 $x_0\in(a,b)$. 由 $f(x)$ 的连续性, 存在 $\delta_0>0$, 使得 $[x_0-\delta_0,x_0+\delta_0]\subset[a,b]$, 且当 $x\in[x_0-\delta_0,x_0+\delta_0]$ 时, 有
\[
  f(x)\geq\frac{f(x_0)}{2}.
\]
因此
\[
  \int_{x_0-\delta_0}^{x_0+\delta_0} f(x)\mathrm dx
  \geq \frac{f(x_0)}{2}\int_{x_0-\delta_0}^{x_0+\delta_0}\mathrm dx
  =\delta_0f(x_0)>0.
\]
这与已知矛盾. (1) 得证。

(2) 用反证法. 当 $f(x)$ 不恒为 $0$ 时, 必存在 $x_0\in(a,b)$, 使得 $f(x_0)>0$. 由 (1) 的证明过程知, 存在 $\delta_0>0$, 使得
\[
  \int_{x_0-\delta_0}^{x_0+\delta_0} f(x)\mathrm dx>0.
\]
但此时有
\[
  \int_a^b f(x)\mathrm dx\geq \int_{x_0-\delta_0}^{x_0+\delta_0} f(x)\mathrm dx>0.
\]
这与已知矛盾. (2) 得证。

\noindent\textbf{例 7.3.3}\quad 设 $0<q\leq p$, 证明 $\ln\dfrac{p}{q}\leq\dfrac{p-q}{q}$。

\noindent\textbf{证明}\quad 当 $p=q$ 时, 结论显然成立. 现在设 $q<p$. 对于 $x\in[q,p]$, 显然有 $\dfrac{1}{x}\leq\dfrac{1}{q}$. 注意到 $\dfrac{1}{x}\in C[q,p]$, 因此 $\dfrac{1}{x}\in R[q,p]$. 再注意到 $\ln x$ 是 $\dfrac{1}{x}$ 的一个原函数, 因此对 $\dfrac{1}{x}\leq\dfrac{1}{q}$ 两边从 $q$ 到 $p$ 积分得
\[
  \ln\frac{p}{q}=\ln x\bigg|_q^p=\int_q^p\frac{1}{x}\mathrm dx
  \leq \int_q^p\frac{1}{q}\mathrm dx=\frac{p-q}{q}.
\]

\section{原函数的存在性与定积分的计算}

\subsection{变限定积分}

在求函数的不定积分时, 我们知道有些初等函数的不定积分不能由初等函数表示. 另外, 当一个函数比较复杂时, 求不定积分也不是一件容易的事. 因此我们很难知道某类函数是否存在原函数. 利用定积分理论, 我们则可以知道对于非常广泛的函数类, 它们总是存在原函数的。

设函数 $f(t)\in R[a,b]$, 则对于 $\forall x\in[a,b]$, $f(t)\in R[a,x]$. 因此
\[
  \Phi(x)=\int_a^x f(t)\mathrm dt
\]
是定义在区间 $[a,b]$ 上的一个函数. 我们称 $\Phi(x)$ 为 $f(t)$ 的变上限定积分. 同样我们也可以在 $[a,b]$ 上定义 $f(t)$ 的变下限定积分
\[
  \Psi(x)=\int_x^b f(t)\mathrm dt\quad(x\in[a,b]).
\]
下面我们将要看到 $\Phi(x),\Psi(x)$ 总比 $f(t)$ 具有更好的性质. 在这里我们只讨论 $\Phi(x)$ 的情况, 对于 $\Psi(x)$ 的情况请读者自己讨论。

\noindent\textbf{定理 7.4.1}\quad 设 $\Phi(x)$ 是函数 $f(t)\in R[a,b]$ 的变上限定积分。

(1) 若 $f(t)\in R[a,b]$, 则 $\Phi(x)\in C[a,b]$;

(2) 若 $f(t)\in C[a,b]$, 则 $\Phi(x)$ 在 $[a,b]$ 上可导, 并且 $\Phi'(x)=f(x)$(在端点处为单侧导数)。

\noindent\textbf{证明}\quad (1) 由 $f(t)\in R[a,b]$, 存在常数 $M>0$, 使得对一切 $t\in[a,b]$, 有 $|f(t)|\leq M$. 任取 $x_0\in[a,b]$, 则当 $x\in[a,b]$ 时, 有
\[
\begin{aligned}
  |\Phi(x)-\Phi(x_0)|
  &=\left|\int_a^x f(t)\mathrm dt-\int_a^{x_0} f(t)\mathrm dt\right|\\
  &=\left|\int_{x_0}^x f(t)\mathrm dt\right|
  \leq \left|\int_{x_0}^x |f(t)|\mathrm dt\right|\\
  &\leq M|x-x_0|.
\end{aligned}
\]
由此推知 $\Phi(x)\in C[a,b]$。

(2) 现设 $f(t)\in C[a,b]$. 下面我们证明 $\Phi(x)$ 在 $x_0\in(a,b)$ 处满足 $\Phi'(x_0)=f(x_0)$(对于 $x_0=a$ 或 $b$ 的情形请读者自证). 由 $f(t)$ 的连续性, 对于 $\forall\varepsilon>0$, $\exists\delta>0$(不妨设 $\delta<\min\{b-x_0,x_0-a\}$), 当 $|t-x_0|<\delta$ 时, 有
\[
  |f(t)-f(x_0)|<\varepsilon.
\]
因此, 当 $0<|x-x_0|<\delta$ 时, 有
\[
\begin{aligned}
  \left|\frac{\Phi(x)-\Phi(x_0)}{x-x_0}-f(x_0)\right|
  &=\left|\frac{\int_a^x f(t)\mathrm dt-\int_a^{x_0} f(t)\mathrm dt}{x-x_0}
  -\frac{\int_{x_0}^x f(x_0)\mathrm dt}{x-x_0}\right|\\
  &=\left|\frac{\int_{x_0}^x [f(t)-f(x_0)]\mathrm dt}{x-x_0}\right|\\
  &\leq \frac{\left|\int_{x_0}^x |f(t)-f(x_0)|\mathrm dt\right|}{|x-x_0|}\\
  &<\varepsilon.
\end{aligned}
\]
这就证明了
\[
  \Phi'(x_0)=\lim_{x\to x_0}\frac{\Phi(x)-\Phi(x_0)}{x-x_0}=f(x_0).
\]
证毕。

由上述定理我们知道, 区间 $[a,b]$ 上的连续函数 $f(x)$ 总存在原函数, 且其变上限定积分 $\Phi(x)=\int_a^x f(t)\mathrm dt(x\in[a,b])$ 即是它的一个原函数. 在以上定理的证明中, 我们实际上证明了更强的结论: 若函数 $f(t)\in R[a,b]$, 并且在 $x_0\in[a,b]$ 处连续, 则其变上限定积分 $\Phi(x)$ 在 $x_0$ 处可导, 并且有 $\Phi'(x_0)=f(x_0)$。

\noindent\textbf{注 1}\quad 若假定函数 $f(x)$ 在区间 $[a,b]$ 上连续, 从定理 7.4.1 可以简单地推导出牛顿--莱布尼茨公式. 事实上, 设 $F(x)$ 是 $f(x)$ 的任一原函数, 则 $F(x)$ 与 $\int_a^x f(t)\mathrm dt$ 同为 $f(x)$ 的原函数, 因此存在常数 $C$, 使得
\[
  F(x)=\int_a^x f(t)\mathrm dt+C.
\]
在上式中, 令 $x=a$, 得 $F(a)=C$; 令 $x=b$, 得
\[
  \int_a^b f(t)\mathrm dt=F(b)-F(a).
\]
再将积分变量换回 $x$ 即得牛顿--莱布尼茨公式。

\noindent\textbf{注 2}\quad 牛顿--莱布尼茨公式是定积分计算的基本方法. 但值得指出的是: 即使一个函数 $f(x)$ 存在原函数 $F(x)$, 有时候 $f(x)$ 也不可积. 例如
\[
  F(x)=
  \begin{cases}
    x^2\sin\dfrac{1}{x^2}, & x\neq0,\\
    0, & x=0,
  \end{cases}
\]
它是
\[
  f(x)=
  \begin{cases}
    2x\sin\dfrac{1}{x^2}-\dfrac{2}{x}\cos\dfrac{1}{x^2}, & x\neq0,\\
    0, & x=0
  \end{cases}
\]
的原函数. 显然 $f(x)$ 在 $x=0$ 的邻域内无界, 因此 $f(x)$ 在包含 $x=0$ 的任何闭区间 $[a,b]$ 上不可积. 对于这类存在原函数的无界函数积分问题, 在下一章中我们还将仔细研究。

下面我们进一步考虑积分上(下)限是函数时变限定积分的求导问题。

若函数 $f(t)\in C[a,b]$, $u=\varphi(x)$ 在区间 $[\alpha,\beta]$ 上可导, 且对于 $\forall x\in[\alpha,\beta]$, $\varphi(x)\in[a,b]$, 则对于 $\forall x\in[\alpha,\beta]$, $\int_a^{\varphi(x)} f(t)\mathrm dt$ 在 $[\alpha,\beta]$ 上有定义, 因此它是 $[\alpha,\beta]$ 上的一个函数. 若令
\[
  G(x)=\int_a^{\varphi(x)} f(t)\mathrm dt,
\]
则 $G(x)$ 可以看成 $\Phi(u)=\int_a^u f(t)\mathrm dt$ 与 $u=\varphi(x)$ 的复合函数. 因此
\[
  G'(x)=\Phi'(\varphi(x))\varphi'(x)=f(\varphi(x))\varphi'(x).
\]
类似地, 我们也可考虑积分下限是函数的情况。

\noindent\textbf{例 7.4.1}\quad 求导数 $\dfrac{\mathrm d}{\mathrm dx}\left(\int_x^{x^2} e^{-t^2}\mathrm dt\right)$。

\noindent\textbf{解}\quad 由于 $e^{-x^2}\in C(-\infty,+\infty)$, 因此它在任何有限闭区间上均可积. 显然有
\[
  \int_x^{x^2} e^{-t^2}\mathrm dt=\int_0^{x^2} e^{-t^2}\mathrm dt-\int_0^x e^{-t^2}\mathrm dt.
\]
令 $\Phi(u)=\int_0^u e^{-t^2}\mathrm dt$, $u=x^2$, 则应用复合函数求导法则得
\[
\begin{aligned}
  \frac{\mathrm d}{\mathrm dx}\left(\int_x^{x^2} e^{-t^2}\mathrm dt\right)
  &=\frac{\mathrm d}{\mathrm dx}\left(\int_0^{x^2} e^{-t^2}\mathrm dt-\int_0^x e^{-t^2}\mathrm dt\right)\\
  &=\left.\frac{\mathrm d(\Phi(u))}{\mathrm du}\right|_{u=x^2}\frac{\mathrm du}{\mathrm dx}-e^{-x^2}\\
  &=2xe^{-x^4}-e^{-x^2}.
\end{aligned}
\]

\noindent\textbf{例 7.4.2}\quad 求极限 $\lim\limits_{x\to0}\dfrac{\int_0^{x^2}\sin t\mathrm dt}{x^4}$。

\noindent\textbf{解}\quad 此极限为 $\dfrac{0}{0}$ 型不定式. 利用洛必达法则得
\[
  \lim_{x\to0}\frac{\int_0^{x^2}\sin t\mathrm dt}{x^4}
  =\lim_{x\to0}\frac{2x\sin x^2}{4x^3}=\frac{1}{2}.
\]

\subsection{定积分的计算}

从牛顿--莱布尼茨公式知, 对于连续函数的定积分, 只要能找到它的一个原函数, 则定积分的问题即可解决. 因此计算不定积分的公式以及求不定积分的方法均可直接应用到定积分的计算中。

\noindent\textbf{例 7.4.3}\quad 求 $I=\int_{-\frac{\pi}{2}}^{\frac{\pi}{2}} f(x)\mathrm dx$, 其中
\[
  f(x)=
  \begin{cases}
    -\sin x, & x\geq0,\\
    x, & x<0.
  \end{cases}
\]

\noindent\textbf{解}\quad 注意到 $f(x)$ 是分段函数, 我们有
\[
\begin{aligned}
  I&=\int_{-\frac{\pi}{2}}^0 x\mathrm dx-\int_0^{\frac{\pi}{2}}\sin x\mathrm dx\\
  &=\left.\frac{x^2}{2}\right|_{-\frac{\pi}{2}}^0+\left.\cos x\right|_0^{\frac{\pi}{2}}
  =-\left(\frac{\pi^2}{8}+1\right).
\end{aligned}
\]
尽管定积分的计算可以先由不定积分找出原函数, 然后计算原函数端点值的差而得到, 但在定积分的计算中, 有时要比不定积分来得简单. 如在不定积分的换元法中必须将所换变量换回原来的变量, 但在定积分的换元法中, 则无此必要。

\noindent\textbf{定理 7.4.2(换元法)}\quad 设函数 $f(x)\in C[a,b]$, $\varphi(t)$ 在区间 $[\alpha,\beta]$ 上具有连续导数, 且 $\varphi(\alpha)=a,\varphi(\beta)=b$, $a\leq\varphi(t)\leq b(\alpha\leq t\leq\beta)$, 则有
\[
  \int_a^b f(x)\mathrm dx=\int_\alpha^\beta f(\varphi(t))\varphi'(t)\mathrm dt.
\]

\noindent\textbf{证明}\quad 一方面, 由于 $f(x)\in C[a,b]$, 因此它在区间 $[a,b]$ 上存在原函数 $F(x)$. 由牛顿--莱布尼茨公式, 我们有
\[
  \int_a^b f(x)\mathrm dx=F(b)-F(a).
\]
另一方面, $F(\varphi(t))$ 是连续函数 $f(\varphi(t))\varphi'(t)$ 在 $[\alpha,\beta]$ 上的一个原函数, 因此也有
\[
\begin{aligned}
  \int_\alpha^\beta f(\varphi(t))\varphi'(t)\mathrm dt
  &=\left.F(\varphi(t))\right|_\alpha^\beta
  =F(\varphi(\beta))-F(\varphi(\alpha))\\
  &=F(b)-F(a).
\end{aligned}
\]
证毕。

\noindent\textbf{注}\quad 若在定理 7.4.2 中, 假定函数 $\varphi(t)$ 在区间 $[\alpha,\beta]$ 上具有连续导数, 且 $\varphi(\alpha)=b,\varphi(\beta)=a,a\leq\varphi(t)\leq b(\alpha\leq t\leq\beta)$, 则有
\[
  \int_a^b f(x)\mathrm dx=\int_\beta^\alpha f(\varphi(t))\varphi'(t)\mathrm dt.
\]
另外, 函数 $\varphi'(t)$ 在 $[\alpha,\beta]$ 上连续的条件可以减弱为 $\varphi'(t)\in R[\alpha,\beta]$. 关于此结果的证明请读者自己给出。

\noindent\textbf{例 7.4.4}\quad 设函数 $f(x)\in C[-a,a](a>0)$, 证明:

(1) 若 $f(x)$ 为偶函数, 则 $\int_{-a}^a f(x)\mathrm dx=2\int_0^a f(x)\mathrm dx$;

(2) 若 $f(x)$ 为奇函数, 则 $\int_{-a}^a f(x)\mathrm dx=0$。

\noindent\textbf{证明}\quad 在 $\int_{-a}^0 f(x)\mathrm dx$ 中令 $x=-t$, 则有
\[
  \int_{-a}^0 f(x)\mathrm dx=-\int_a^0 f(-t)\mathrm dt=\int_0^a f(-x)\mathrm dx.
\]
当 $f(x)$ 是偶函数时, 我们有 $f(-x)=f(x)(x\in[-a,a])$. 因此由定积分的性质有
\[
\begin{aligned}
  \int_{-a}^a f(x)\mathrm dx
  &=\int_{-a}^0 f(x)\mathrm dx+\int_0^a f(x)\mathrm dx\\
  &=\int_0^a f(-x)\mathrm dx+\int_0^a f(x)\mathrm dx
  =2\int_0^a f(x)\mathrm dx.
\end{aligned}
\]
(1) 得证。

当 $f(x)$ 是奇函数时, 我们有 $f(-x)=-f(x)(x\in[-a,a])$. 因此由定积分的性质有
\[
\begin{aligned}
  \int_{-a}^a f(x)\mathrm dx
  &=\int_{-a}^0 f(x)\mathrm dx+\int_0^a f(x)\mathrm dx\\
  &=\int_0^a f(-x)\mathrm dx+\int_0^a f(x)\mathrm dx\\
  &=-\int_0^a f(x)\mathrm dx+\int_0^a f(x)\mathrm dx=0.
\end{aligned}
\]
(2) 得证。

\noindent\textbf{例 7.4.5}\quad 求定积分 $\int_0^a \sqrt{a^2-x^2}\mathrm dx$, $a>0$。

\noindent\textbf{解}\quad 令 $x=a\sin t\left(t\in\left[0,\dfrac{\pi}{2}\right]\right)$, 则 $x=0\Rightarrow t=0$, $x=a\Rightarrow t=\dfrac{\pi}{2}$. 由此得
\[
\begin{aligned}
  \int_0^a \sqrt{a^2-x^2}\mathrm dx
  &=\int_0^{\frac{\pi}{2}} a\sqrt{a^2\cos^2 t}\cos t\mathrm dt\\
  &=a^2\int_0^{\frac{\pi}{2}}\cos^2 t\mathrm dt
  =a^2\int_0^{\frac{\pi}{2}}\frac{1+\cos 2t}{2}\mathrm dt\\
  &=\frac{a^2}{2}\left.\left(t+\frac{1}{2}\sin 2t\right)\right|_0^{\frac{\pi}{2}}
  =\frac{\pi}{4}a^2.
\end{aligned}
\]

\noindent\textbf{例 7.4.6}\quad 设 $f(x)$ 是以 $T(T>0)$ 为周期的连续函数, 证明: 对于任意的 $a$, 有
\[
  \int_a^{a+T} f(x)\mathrm dx=\int_0^T f(x)\mathrm dx.
\]

\noindent\textbf{证明}\quad 由定积分的性质, 我们有
\begin{equation}
  \int_a^{a+T} f(x)\mathrm dx
  =\int_a^0 f(x)\mathrm dx+\int_0^T f(x)\mathrm dx+\int_T^{a+T} f(x)\mathrm dx.
  \tag{7.3.7}
\end{equation}
令 $x=t+T$, 则有
\[
  \int_T^{a+T} f(x)\mathrm dx
  =\int_0^a f(t+T)\mathrm dt
  =\int_0^a f(t)\mathrm dt
  =-\int_a^0 f(t)\mathrm dt.
\]
将此等式中的积分变量 $t$ 换回 $x$ 后, 我们有
\[
  \int_T^{a+T} f(x)\mathrm dx=-\int_a^0 f(x)\mathrm dx.
\]
将此等式代入 (7.3.7) 即得所证。

\noindent\textbf{例 7.4.7}\quad 求定积分 $I=\int_0^\pi \dfrac{x\sin x}{1+\cos^2 x}\mathrm dx$。

\noindent\textbf{解}\quad 由定积分的性质, 我们有
\[
  I=\int_0^{\frac{\pi}{2}}\frac{x\sin x}{1+\cos^2 x}\mathrm dx
  +\int_{\frac{\pi}{2}}^\pi \frac{x\sin x}{1+\cos^2 x}\mathrm dx.
\]
在 $\int_{\frac{\pi}{2}}^\pi \dfrac{x\sin x}{1+\cos^2 x}\mathrm dx$ 中令 $x=\pi-t$, 则 $x=\dfrac{\pi}{2}\Rightarrow t=\dfrac{\pi}{2},x=\pi\Rightarrow t=0$, 从而有
\[
\begin{aligned}
  \int_{\frac{\pi}{2}}^\pi \frac{x\sin x}{1+\cos^2 x}\mathrm dx
  &=-\int_{\frac{\pi}{2}}^0
  \frac{(\pi-t)\sin(\pi-t)}{1+\cos^2(\pi-t)}\mathrm dt\\
  &=\pi\int_0^{\frac{\pi}{2}}\frac{\sin x}{1+\cos^2 x}\mathrm dx
  -\int_0^{\frac{\pi}{2}}\frac{x\sin x}{1+\cos^2 x}\mathrm dx.
\end{aligned}
\]
最后我们得到
\[
  I=\pi\int_0^{\frac{\pi}{2}}\frac{\sin x}{1+\cos^2 x}\mathrm dx
  =-\left.\pi\arctan(\cos x)\right|_0^{\frac{\pi}{2}}
  =\frac{\pi^2}{4}.
\]

\noindent\textbf{例 7.4.8}\quad 设 $x_1,x_2\in(-1,1)$, 定义
\[
  d(x_1,x_2)=\left|\int_{x_1}^{x_2}\frac{\mathrm dx}{1-x^2}\right|.
\]
试证明以下结论:

(1) 对任何 $x_1,x_2\in(-1,1)$, 有 $d(x_1,x_2)\geq0$, 且
\[
  d(x_1,x_2)=0\Longleftrightarrow x_1=x_2;
\]

(2) $d(x_1,x_2)=d(x_2,x_1)$;

(3) 对任何 $x_1,x_2,x_3\in(-1,1)$, 成立
\[
  d(x_1,x_3)\leq d(x_1,x_2)+d(x_2,x_3);
\]

(4) 设 $x_0,x_1,x_2\in(-1,1)$, $w_j=\dfrac{x_j-x_0}{1-x_0x_j}(j=1,2)$, 则 $w_1,w_2\in(-1,1)$, 并且 $d(w_1,w_2)=d(x_1,x_2)$;

(5) 设 $x_0\in(-1,1)$, 则 $\lim\limits_{x\to\pm1}d(x_0,x)=+\infty$。

\noindent\textbf{证明}\quad (1) 当 $x_1=x_2$ 时, 显然 $d(x_1,x_2)=0$. 当 $x_1\neq x_2$ 时, 有
\[
  d(x_1,x_2)
  =\left|\int_{x_1}^{x_2}\frac{\mathrm dx}{1-x^2}\right|
  =\left|\frac{1}{2}\left(\ln\frac{1+x_2}{1-x_2}-\ln\frac{1+x_1}{1-x_1}\right)\right|>0.
\]
因此若 $d(x_1,x_2)=0$, 必有 $x_1=x_2$, 从而 (1) 成立。

(2) 由于 $\int_{x_1}^{x_2}\dfrac{\mathrm dx}{1-x^2}=-\int_{x_2}^{x_1}\dfrac{\mathrm dx}{1-x^2}$, 两边取绝对值即得 (2)。

(3) 当以 $x_1,x_3$ 为端点的开区间包含在以 $x_1,x_2$ 或 $x_2,x_3$ 为端点的开区间时, 显然结论成立. 现不妨设 $x_1<x_2<x_3$, 则有
\[
\begin{aligned}
  d(x_1,x_3)
  &=\int_{x_1}^{x_3}\frac{\mathrm dx}{1-x^2}
  =\int_{x_1}^{x_2}\frac{\mathrm dx}{1-x^2}
  +\int_{x_2}^{x_3}\frac{\mathrm dx}{1-x^2}\\
  &=d(x_1,x_2)+d(x_2,x_3).
\end{aligned}
\]

(4) 当 $x_0,x_1\in(-1,1)$ 时, 我们有
\[
  0<(1-x_0^2)(1-x_1^2)=1-(x_0^2+x_1^2)+x_0^2x_1^2,
\]
从而有
\[
  (x_1-x_0)^2=x_1^2+x_0^2-2x_0x_1<1+x_0^2x_1^2-2x_0x_1=(1-x_0x_1)^2.
\]
因此有 $w_1\in(-1,1)$. 同理有 $w_2\in(-1,1)$。

当 $x_0\in(-1,1)$ 时, $w(x)=\dfrac{x-x_0}{1-x_0x}$ 在 $x\in(-1,1)$ 时有定义且有
\[
  w'(x)=\frac{1-x_0^2}{(1-x_0x)^2}>0,
\]
因此 $w(x)$ 在 $(-1,1)$ 内严格单调上升. 注意到我们曾经证明过等式
\[
  \frac{\mathrm dw}{1-w^2}=\frac{\mathrm dx}{1-x^2},
\]
因此由定积分的换元法有
\[
  d(x_1,x_2)=\left|\int_{x_1}^{x_2}\frac{\mathrm dx}{1-x^2}\right|
  =\left|\int_{w_1}^{w_2}\frac{\mathrm dw}{1-w^2}\right|
  =d(w_1,w_2).
\]

(5) 直接计算可得
\[
\begin{aligned}
  \lim_{x\to\pm1}d(x_0,x)
  &=\lim_{x\to\pm1}\left|\int_{x_0}^x\frac{\mathrm dx}{1-x^2}\right|\\
  &=\lim_{x\to\pm1}\frac{1}{2}
  \left|\ln\frac{1+x}{1-x}-\ln\frac{1+x_0}{1-x_0}\right|\\
  &=+\infty.
\end{aligned}
\]

\noindent\textbf{注}\quad 若在区间 $(-1,1)$ 内如上定义两点间距离 $d$, 则称该距离为庞加莱 (Poincaré) 距离或双曲距离. 容易看出, 在该距离下, $(-1,1)$ 内的柯西序列都是收敛的, 并且其极限是 $(-1,1)$ 内的点. 因此在该距离下, $(-1,1)$ 是完备的。

\noindent\textbf{定理 7.4.3(分部积分法)}\quad 设函数 $u(x),v(x)$ 在区间 $[a,b]$ 上可导, 并且 $u'(x),v'(x)\in R[a,b]$, 则
\[
  \int_a^b u(x)v'(x)\mathrm dx
  =\left.u(x)v(x)\right|_a^b-\int_a^b u'(x)v(x)\mathrm dx.
\]

\noindent\textbf{证明}\quad 在恒等式 $[u(x)v(x)]'=u'(x)v(x)+u(x)v'(x)$ 中, 我们知等式两边的函数在 $[a,b]$ 上均可积, 因此两边从 $a$ 到 $b$ 对 $x$ 积分并利用牛顿--莱布尼茨公式得
\[
  \left.u(x)v(x)\right|_a^b
  =\int_a^b [u(x)v(x)]'\mathrm dx
  =\int_a^b u'(x)v(x)\mathrm dx+\int_a^b u(x)v'(x)\mathrm dx.
\]
将上述等式整理即得所证。

\noindent\textbf{例 7.4.9}\quad 计算定积分 $I=\int_0^1 x\ln(1+x)\mathrm dx$。

\noindent\textbf{解}\quad
\[
\begin{aligned}
  I&=\int_0^1 \ln(1+x)\mathrm d\left(\frac{x^2}{2}\right)
  =\left.\frac{x^2}{2}\ln(1+x)\right|_0^1
  -\int_0^1 \frac{x^2}{2}\mathrm d(\ln(1+x))\\
  &=\frac{1}{2}\ln2-\frac{1}{2}\int_0^1\left(x-1+\frac{1}{1+x}\right)\mathrm dx\\
  &=\frac{1}{2}\ln2-\frac{1}{2}\left.\left[\frac{x^2}{2}-x+\ln(1+x)\right]\right|_0^1
  =\frac{1}{4}.
\end{aligned}
\]

\noindent\textbf{例 7.4.10}\quad 证明: 对任意的 $n\in\mathbb N$, 有
\[
  \int_0^{\frac{\pi}{2}}\sin^n x\mathrm dx
  =\int_0^{\frac{\pi}{2}}\cos^n x\mathrm dx\triangleq I_n,
\]
并求 $I_n$ 的值。

\noindent\textbf{证明}\quad 对 $\int_0^{\frac{\pi}{2}}\sin^n x\mathrm dx$ 作变换 $x=\dfrac{\pi}{2}-t$, 得
\[
  \int_0^{\frac{\pi}{2}}\sin^n x\mathrm dx
  =\int_{\frac{\pi}{2}}^0 \cos^n t\,\mathrm d\left(\frac{\pi}{2}-t\right)
  =-\int_{\frac{\pi}{2}}^0 \cos^n t\mathrm dt
  =\int_0^{\frac{\pi}{2}}\cos^n x\mathrm dx.
\]
下面我们来求 $I_n$ 的递推公式. 对于 $n=0,1$, 我们有
\[
  I_0=\int_0^{\frac{\pi}{2}}\mathrm dx=\frac{\pi}{2},\qquad
  I_1=\int_0^{\frac{\pi}{2}}\sin x\mathrm dx=1.
\]
当 $n>1$ 时, 有
\[
\begin{aligned}
  I_n&=\int_0^{\frac{\pi}{2}}\sin^{n-1}x\,\mathrm d(-\cos x)\\
  &=-\left.\sin^{n-1}x\cos x\right|_0^{\frac{\pi}{2}}
  +\int_0^{\frac{\pi}{2}}\cos x\,\mathrm d(\sin^{n-1}x)\\
  &=(n-1)\int_0^{\frac{\pi}{2}}\sin^{n-2}x(1-\sin^2x)\mathrm dx\\
  &=(n-1)I_{n-2}-(n-1)I_n,
\end{aligned}
\]
即 $I_n=\dfrac{n-1}{n}I_{n-2}$. 因此得
\[
  I_{2k}=\frac{(2k-1)!!}{(2k)!!}\cdot\frac{\pi}{2},\qquad
  I_{2k+1}=\frac{(2k)!!}{(2k+1)!!}\quad(k=1,2,\cdots).
\]
作为上述定积分的应用, 可以推出下面的瓦利斯 (Wallis) 公式:
\[
  \lim_{n\to\infty}\left[\frac{(2n)!!}{(2n-1)!!}\right]^2\frac{1}{2n+1}=\frac{\pi}{2}.
\]
事实上, 当 $x\in\left(0,\dfrac{\pi}{2}\right)$ 时, 对任意 $n\in\mathbb N$, 有
\[
  \sin^{2n+1}x<\sin^{2n}x<\sin^{2n-1}x,
\]
从而有
\[
  \int_0^{\frac{\pi}{2}}\sin^{2n+1}x\mathrm dx
  <\int_0^{\frac{\pi}{2}}\sin^{2n}x\mathrm dx
  <\int_0^{\frac{\pi}{2}}\sin^{2n-1}x\mathrm dx.
\]
利用例 7.4.10 有
\[
  \frac{(2n)!!}{(2n+1)!!}
  <\frac{(2n-1)!!}{(2n)!!}\cdot\frac{\pi}{2}
  <\frac{(2n-2)!!}{(2n-1)!!}.
\]
从上式容易推出
\[
  J_n\triangleq\left[\frac{(2n)!!}{(2n-1)!!}\right]^2\frac{1}{2n+1}
  <\frac{\pi}{2}
  <\left[\frac{(2n)!!}{(2n-1)!!}\right]^2\frac{1}{2n}\triangleq J'_n.
\]
由于
\[
\begin{aligned}
  0&\leq\lim_{n\to\infty}(J'_n-J_n)\\
  &=\lim_{n\to\infty}
  \left[\frac{(2n)!!}{(2n-1)!!}\right]^2\frac{1}{2n(2n+1)}\\
  &\leq\lim_{n\to\infty}\frac{1}{2n}\cdot\frac{\pi}{2}=0,
\end{aligned}
\]
因此瓦利斯公式成立。

\section{定积分中值定理}

在定积分的理论中, 我们最后介绍定积分的两个重要性质。

\subsection{定积分第一中值定理}

设函数 $f(x)$ 在区间 $[a,b]$ 上连续, 由定积分的几何意义知, 必存在一点 $\xi\in[a,b]$, 使得
\[
  \int_a^b f(x)\mathrm dx=f(\xi)\int_a^b\mathrm dx=f(\xi)(b-a).
\]
这里 $\int_a^b f(x)\mathrm dx$ 可以看成函数 $f(x)g(x)=f(x)\cdot1$ 在 $[a,b]$ 上的定积分, 其中 $g(x)\equiv1$. 一个自然的问题是: 函数 $g(x)\equiv1$ 能否推广成一般的函数? 定积分第一中值定理将回答这个问题。

\noindent\textbf{定理 7.5.1(定积分第一中值定理)}\quad 设函数 $f(x)\in C[a,b],g(x)\in R[a,b]$ 且在区间 $[a,b]$ 上不变号, 则存在 $\xi\in[a,b]$, 使得
\[
  \int_a^b f(x)g(x)\mathrm dx=f(\xi)\int_a^b g(x)\mathrm dx.
\]
特别地, 当 $g(x)\equiv1$ 时, 有 $\int_a^b f(x)\mathrm dx=f(\xi)(b-a)$。

\noindent\textbf{证明}\quad 不妨设 $g(x)\geq0$, 并令 $m,M$ 分别是函数 $f(x)$ 在区间 $[a,b]$ 上的最小值和最大值, 则对于 $\forall x\in[a,b]$, 有
\[
  mg(x)\leq f(x)g(x)\leq Mg(x).
\]
因此有
\[
  m\int_a^b g(x)\mathrm dx
  \leq\int_a^b f(x)g(x)\mathrm dx
  \leq M\int_a^b g(x)\mathrm dx.
\]
若 $\int_a^b g(x)\mathrm dx=0$, 由上式得 $\int_a^b f(x)g(x)\mathrm dx=0$, 此时任取 $\xi\in[a,b]$ 即可. 当 $\int_a^b g(x)\mathrm dx\neq0$ 时, 则有
\[
  m\leq
  \frac{\int_a^b f(x)g(x)\mathrm dx}{\int_a^b g(x)\mathrm dx}
  \leq M.
\]
由连续函数的介值定理, 必存在 $\xi\in[a,b]$, 使得
\[
  f(\xi)=
  \frac{\int_a^b f(x)g(x)\mathrm dx}{\int_a^b g(x)\mathrm dx}.
\]
证毕。

当 $g(x)\equiv1$ 且 $f(x)\geq0$ 时, 定积分第一中值定理具有鲜明的几何意义, 即由直线 $x=a,x=b,y=0$ 和曲线 $y=f(x)$ 所围的曲边梯形的面积必与由直线 $x=a,x=b,y=0,y=f(\xi)(\xi\in[a,b])$ 围成的矩形面积相等 (见图 7.5.1)。

\noindent\textbf{注 1}\quad 在定积分第一中值定理中, 若只假定 $f(x)\in R[a,b]$, 令
\[
  m=\inf_{x\in[a,b]}\{f(x)\},\qquad
  M=\sup_{x\in[a,b]}\{f(x)\},
\]
则用同样的方法可以证明存在 $\mu\in[m,M]$, 使得
\[
  \int_a^b f(x)g(x)\mathrm dx=\mu\int_a^b g(x)\mathrm dx.
\]

\noindent\textbf{注 2}\quad 若在定积分第一中值定理中假定 $f(x)$ 与 $g(x)$ 均是连续函数, 且 $g(x)$ 在 $[a,b]$ 上不变号, 则可以证明必存在 $\xi\in(a,b)$, 使得
\[
  \int_a^b f(x)g(x)\mathrm dx=f(\xi)\int_a^b g(x)\mathrm dx
\]
(见本章习题)。

\noindent\textbf{例 7.5.1}\quad 设函数 $f(x)\in C[0,1]$, 证明
\[
  \lim_{n\to\infty}\int_0^1\frac{nf(x)}{1+n^2x^2}\mathrm dx=\frac{\pi}{2}f(0).
\]

\noindent\textbf{证明}\quad 由于函数 $f(x)\in C[0,1]$, 从而存在 $M>0$, 使得对于 $\forall x\in[0,1]$, 有 $|f(x)|\leq M$. 由定积分的性质, 对于 $\forall n\in\mathbb N$, 我们有
\[
  \int_0^1\frac{nf(x)}{1+n^2x^2}\mathrm dx
  =\int_0^{n^{-\frac13}}\frac{nf(x)}{1+n^2x^2}\mathrm dx
  +\int_{n^{-\frac13}}^1\frac{nf(x)}{1+n^2x^2}\mathrm dx.
\]
下面我们分别来估计上面等式右边中的两个定积分. 对于第二个定积分我们有
\[
  \left|\int_{n^{-\frac13}}^1\frac{nf(x)}{1+n^2x^2}\mathrm dx\right|
  \leq\int_{n^{-\frac13}}^1\left|\frac{nf(x)}{1+n^2x^2}\right|\mathrm dx
  \leq\frac{nM}{1+n^{1+\frac13}}\to0\quad(n\to\infty).
\]
而在第一个定积分 $\int_0^{n^{-\frac13}}\dfrac{nf(x)}{1+n^2x^2}\mathrm dx$ 中, $\dfrac{n}{1+n^2x^2}$ 在 $[0,1]$ 上不变号, 因此存在 $\xi\in[0,n^{-\frac13}]$, 使得
\[
\begin{aligned}
  \int_0^{n^{-\frac13}}\frac{nf(x)}{1+n^2x^2}\mathrm dx
  &=f(\xi)\int_0^{n^{-\frac13}}\frac{n}{1+n^2x^2}\mathrm dx\\
  &=\left.f(\xi)\arctan nx\right|_0^{n^{-\frac13}}
  =f(\xi)\arctan n^{\frac23}.
\end{aligned}
\]
当 $n\to\infty$ 时, 由于 $\xi\in[0,n^{-\frac13}]$, 有 $\xi\to0$, 而 $\arctan n^{\frac23}\to\dfrac{\pi}{2}$, 因此
\[
  \int_0^{n^{-\frac13}}\frac{nf(x)}{1+n^2x^2}\mathrm dx\to\frac{\pi}{2}f(0).
\]
最终我们有
\[
\begin{aligned}
  \lim_{n\to\infty}\int_0^1\frac{nf(x)}{1+n^2x^2}\mathrm dx
  &=\lim_{n\to\infty}\int_0^{n^{-\frac13}}\frac{nf(x)}{1+n^2x^2}\mathrm dx
  +\lim_{n\to\infty}\int_{n^{-\frac13}}^1\frac{nf(x)}{1+n^2x^2}\mathrm dx\\
  &=\frac{\pi}{2}f(0).
\end{aligned}
\]

利用定积分的分部积分法我们可以得到带积分余项的泰勒公式; 再利用定积分第一中值定理, 我们还可以从带积分余项的泰勒公式推出带拉格朗日余项和带柯西余项的泰勒公式。

\noindent\textbf{定理 7.5.2}\quad 设函数 $f(x)$ 在 $x_0$ 的邻域 $U(x_0,h)(h>0)$ 内具有 $n+1$ 阶连续导数 $f^{(n+1)}(x)$, 则对 $\forall x\in U(x_0,h)$, 有
\begin{equation}
\begin{aligned}
  f(x)&=f(x_0)+f'(x_0)(x-x_0)+\frac{f''(x_0)}{2!}(x-x_0)^2\\
  &\quad+\cdots+\frac{f^{(n)}(x_0)}{n!}(x-x_0)^n+R_n(x),
\end{aligned}
\tag{7.5.1}
\end{equation}
其中
\[
  R_n(x)=\frac{1}{n!}\int_{x_0}^x f^{(n+1)}(t)(x-t)^n\mathrm dt.
\]

\noindent\textbf{证明}\quad 由牛顿--莱布尼茨公式有
\[
  f(x)=f(x_0)+\int_{x_0}^x f'(t)\mathrm dt
  =f(x_0)+\int_{x_0}^x (x-t)^0f'(t)\mathrm dt.
\]
利用分部积分我们有
\[
\begin{aligned}
  \int_{x_0}^x (x-t)^0f'(t)\mathrm dt
  &=-\int_{x_0}^x f'(t)\mathrm d(x-t)\\
  &=-\left.f'(t)(x-t)\right|_{x_0}^x+\int_{x_0}^x (x-t)f''(t)\mathrm dt\\
  &=f'(x_0)(x-x_0)+\int_{x_0}^x (x-t)f''(t)\mathrm dt\\
  &=f'(x_0)(x-x_0)-\frac{1}{2}\int_{x_0}^x f''(t)\mathrm d(x-t)^2\\
  &=f'(x_0)(x-x_0)-\left.\frac{f''(t)}{2}(x-t)^2\right|_{x_0}^x
  +\frac{1}{2}\int_{x_0}^x (x-t)^2f'''(t)\mathrm dt\\
  &=f'(x_0)(x-x_0)+\frac{f''(x_0)}{2}(x-x_0)^2
  +\frac{1}{2}\int_{x_0}^x (x-t)^2f'''(t)\mathrm dt.
\end{aligned}
\]
再将上式中的积分继续分部积分, 逐次类推即可得到所要结论. 证毕。

\noindent\textbf{注 1}\quad 定理 7.5.2 中的带余项
\[
  R_n(x)=\frac{1}{n!}\int_{x_0}^x f^{(n+1)}(t)(x-t)^n\mathrm dt
\]
的公式 (7.5.1) 称为带积分余项的泰勒公式, 该余项也称为积分余项. 在定理 7.5.2 的条件下, 我们可以推出带拉格朗日余项的泰勒公式. 事实上, 在 $x_0$ 与 $x$ 之间, $g(t)=(x-t)^n$ 是关于 $t$ 的连续函数并且是不变号的, 而 $f^{(n+1)}(t)$ 是连续的, 因此, 由定积分第一中值定理的注 2, 在 $x$ 与 $x_0$ 之间存在 $\xi$, 使得下式成立:
\[
\begin{aligned}
  R_n(x)&=\frac{1}{n!}\int_{x_0}^x f^{(n+1)}(t)(x-t)^n\mathrm dt\\
  &=\frac{f^{(n+1)}(\xi)}{n!}\int_{x_0}^x (x-t)^n\mathrm dt\\
  &=\frac{f^{(n+1)}(\xi)}{(n+1)!}(x-x_0)^{n+1}.
\end{aligned}
\]
这就是在泰勒公式中的拉格朗日余项。

\noindent\textbf{注 2}\quad 从定理 7.5.2 中我们还可以推出带以下柯西余项的泰勒公式, 即在定理 7.5.2 的条件下, 我们有
\begin{equation}
  R_n(x)=\frac{1}{n!}f^{(n+1)}[x_0+\theta(x-x_0)](1-\theta)^n(x-x_0)^{n+1},
  \tag{7.5.2}
\end{equation}
其中 $0<\theta<1$. 事实上, 对积分余项应用定积分第一中值定理的注 2 知, 在 $x$ 与 $x_0$ 之间存在 $\xi$, 使得下式成立:
\[
  R_n(x)=\frac{1}{n!}\int_{x_0}^x f^{(n+1)}(t)(x-t)^n\mathrm dt
  =\frac{1}{n!}f^{(n+1)}(\xi)(x-\xi)^n(x-x_0).
\]
注意到 $x-\xi=(x-x_0)-\theta(x-x_0)=(1-\theta)(x-x_0)$, 其中 $0<\theta<1$. 将此代入上式整理即得 (7.5.2)。

形如 (7.5.2) 的余项称为柯西余项, 带此余项的泰勒公式称为带柯西余项的泰勒公式。

\subsection{定积分第二中值定理}

设 $f(x)$ 是区间 $[a,b]$ 上单调上升的函数且 $f(x)\geq0$, 则 $\int_a^b f(x)\mathrm dx\geq0$, 从而有
\[
  f(b)(b-b)\leq\int_a^b f(x)\mathrm dx\leq f(b)(b-a).
\]
于是存在 $\xi\in[a,b]$, 使得
\[
  \int_a^b f(x)\mathrm dx=f(b)(b-\xi)=f(b)\int_\xi^b\mathrm dx.
\]
因此人们不禁要问: 对区间 $[a,b]$ 上的可积函数 $g(x)$, 是否也存在 $\xi\in[a,b]$, 使得 $\int_a^b f(x)g(x)\mathrm dx=f(b)\int_\xi^b g(x)\mathrm dx$? 定积分第二中值定理将回答这个问题, 该定理在后面章节中将多次用到。

\noindent\textbf{定理 7.5.3(定积分第二中值定理)}\quad 设函数 $g(x)\in R[a,b]$。

(1) 若函数 $f(x)$ 在区间 $[a,b]$ 上单调上升, 且对于 $\forall x\in[a,b]$, 有 $f(x)\geq0$, 则存在 $\xi_1\in[a,b]$, 使得
\[
  \int_a^b f(x)g(x)\mathrm dx=f(b)\int_{\xi_1}^b g(x)\mathrm dx;
\]

(2) 若函数 $f(x)$ 在区间 $[a,b]$ 上单调下降, 且对于 $\forall x\in[a,b]$, 有 $f(x)\geq0$, 则存在 $\xi_2\in[a,b]$, 使得
\[
  \int_a^b f(x)g(x)\mathrm dx=f(a)\int_a^{\xi_2} g(x)\mathrm dx;
\]

(3) 若函数 $f(x)$ 在区间 $[a,b]$ 上单调, 则存在 $\xi\in[a,b]$, 使得
\[
  \int_a^b f(x)g(x)\mathrm dx
  =f(a)\int_a^\xi g(x)\mathrm dx+f(b)\int_\xi^b g(x)\mathrm dx.
\]

由于定理 7.5.3 的条件比较弱, 从而它的证明也比较复杂. 为了看出定理 7.5.3 的证明思路, 我们的方法是先尝试证明该定理在一些较强条件下的相同的结果; 然后对其证明过程加以考察, 分析出证明的关键所在, 从而找到在较弱条件下的证明思路. 以 (1) 为例, 若在原来假定下, 再假设 $g(x)\in C[a,b],f(x)\in C^1[a,b]$, 则有 $f'(x)\geq0,x\in[a,b]$, 且 $f(a)\geq0$. 在这种情况下, 我们可以用下面简单方法来证明 (1)。

令 $G(x)=\int_x^b g(t)\mathrm dt$, 则 $G(x)$ 在 $[a,b]$ 上连续可导. 设 $m,M$ 分别为 $G(x)$ 在 $[a,b]$ 上的最小值和最大值. 由 $G'(x)=-g(x),G(b)=0$, 得
\[
\begin{aligned}
  \int_a^b f(x)g(x)\mathrm dx
  &=-\int_a^b f(x)\mathrm dG(x)\\
  &=-\left.f(x)G(x)\right|_a^b+\int_a^b G(x)f'(x)\mathrm dx\\
  &=f(a)G(a)+\int_a^b G(x)f'(x)\mathrm dx.
\end{aligned}
\]
再由
\[
  m[f(b)-f(a)]\leq\int_a^b G(x)f'(x)\mathrm dx\leq M[f(b)-f(a)],
\]
我们推出
\[
  f(a)G(a)+\int_a^b G(x)f'(x)\mathrm dx
  \leq f(a)M+M[f(b)-f(a)]=Mf(b)
\]
和
\[
  f(a)G(a)+\int_a^b G(x)f'(x)\mathrm dx
  \geq f(a)m+m[f(b)-f(a)]=mf(b).
\]
因此有
\[
  mf(b)\leq\int_a^b f(x)g(x)\mathrm dx\leq Mf(b).
\]
若 $f(b)=0$, 则 $f(x)\equiv0$, 结论显然成立. 现设 $f(b)>0$, 则有
\[
  m\leq\frac{\int_a^b f(x)g(x)\mathrm dx}{f(b)}\leq M.
\]
由连续函数的介值定理, 存在 $\xi_1\in[a,b]$, 使得
\[
  G(\xi_1)=\int_{\xi_1}^b g(x)\mathrm dx
  =\frac{1}{f(b)}\int_a^b f(x)g(x)\mathrm dx.
\]
因此 (1) 成立。

从上述讨论可以看出, 若要证明定理 7.5.3 中的 (1), 我们仍然可以设 $G(x)=\int_x^b g(t)\mathrm dt$. 由 $g(x)$ 的可积性知 $G(x)$ 在区间 $[a,b]$ 上连续. 如果在这种情形下还是能证明不等式
\[
  mf(b)\leq\int_a^b f(x)g(x)\mathrm dx\leq Mf(b)
\]
成立, 其中 $m,M$ 分别为 $G(x)$ 在 $[a,b]$ 上的最小值和最大值, 我们就可以证明 (1) 成立. 但现在我们的假定只是 $f(x)$ 与 $g(x)$ 在 $[a,b]$ 上可积, 因此我们只能从定积分的定义出发来证明该结论. 在定理 7.5.3 的证明过程中, 还需要用到下面的阿贝尔变换. 为此, 我们将其以引理的形式给出。

\noindent\textbf{引理 7.5.4(阿贝尔变换)}\quad 设 $\alpha_1,\alpha_2,\cdots,\alpha_n$ 和 $\beta_1,\beta_2,\cdots,\beta_n$ 为两组数, 记 $B_k=\sum_{i=1}^k\beta_i(k=1,2,\cdots,n)$, 则有
\[
  \sum_{i=1}^n\alpha_i\beta_i
  =\sum_{i=1}^{n-1}(\alpha_i-\alpha_{i+1})B_i+\alpha_nB_n.
\]

\noindent\textbf{证明}\quad 由 $\beta_i=B_i-B_{i-1}(i=2,3,\cdots,n)$, 得
\[
\begin{aligned}
  \sum_{i=1}^n\alpha_i\beta_i
  &=\sum_{i=2}^n\alpha_i(B_i-B_{i-1})+\alpha_1B_1\\
  &=B_1(\alpha_1-\alpha_2)+B_2(\alpha_2-\alpha_3)+\cdots\\
  &\quad+B_{n-1}(\alpha_{n-1}-\alpha_n)+B_n\alpha_n.
\end{aligned}
\]
证毕。

\noindent\textbf{定理 7.5.2 的证明}\quad 我们只证 (2) 和 (3), (1) 的证明留给读者. 为了证明 (2), 我们令 $h(x)=\int_a^x g(t)\mathrm dt$. 由变限定积分的性质知 $h(x)\in C[a,b]$. 记
\[
  m=\min_{x\in[a,b]}\{h(x)\},\qquad
  M=\max_{x\in[a,b]}\{h(x)\}.
\]
由于函数 $g(x)\in R[a,b]$, 存在 $0<M_1<+\infty$, 使得对于 $\forall x\in[a,b]$, 有 $|g(x)|\leq M_1$。

设 $\Delta:a=x_0<x_1<\cdots<x_n=b$ 为 $[a,b]$ 的任一个分割, 则有
\[
  \int_a^b f(x)g(x)\mathrm dx
  =\lim_{\lambda(\Delta)\to0}\sum_{i=1}^n\int_{x_{i-1}}^{x_i} f(x)g(x)\mathrm dx.
\]
由于
\[
  \left|\sum_{i=1}^n\int_{x_{i-1}}^{x_i}[f(x)-f(x_{i-1})]g(x)\mathrm dx\right|
  \leq M_1\sum_{i=1}^n\omega_i\Delta x_i\to0\quad(\lambda(\Delta)\to0),
\]
其中 $\omega_i$ 为 $f(x)$ 在 $[x_{i-1},x_i](i=1,2,\cdots,n)$ 上的振幅, 我们有
\[
  \int_a^b f(x)g(x)\mathrm dx
  =\lim_{\lambda(\Delta)\to0}\sum_{i=1}^n f(x_{i-1})\int_{x_{i-1}}^{x_i}g(x)\mathrm dx.
\]
因此我们只要对
\[
  \sum_{i=1}^n f(x_{i-1})\int_{x_{i-1}}^{x_i}g(x)\mathrm dx
  =\sum_{i=1}^n f(x_{i-1})[h(x_i)-h(x_{i-1})]
\]
进行估计即可。

利用阿贝尔变换 (引理 7.5.4), 我们有
\[
\sum_{i=1}^n f(x_{i-1})[h(x_i)-h(x_{i-1})]
=\sum_{i=1}^{n-1}[f(x_{i-1})-f(x_i)]h(x_i)+f(x_{n-1})h(b).
\]
由于 $f(x)$ 在 $[a,b]$ 上是单调下降的, 我们有 $f(x_{i-1})-f(x_i)\geq0(i=1,2,\cdots,n-1)$, 再注意到 $f(b)\geq0$ 和 $m\leq h(x)\leq M$, 可以推出
\[
\begin{aligned}
  mf(a)&=\sum_{i=1}^{n-1}[f(x_{i-1})-f(x_i)]m+f(x_{n-1})m\\
  &\leq\sum_{i=1}^{n-1}[f(x_{i-1})-f(x_i)]h(x_i)+f(x_{n-1})h(b)\\
  &\leq\sum_{i=1}^{n-1}[f(x_{i-1})-f(x_i)]M+f(x_{n-1})M\\
  &=Mf(a).
\end{aligned}
\]
由此得到
\[
  mf(a)\leq\int_a^b f(x)g(x)\mathrm dx
  =\lim_{\lambda(\Delta)\to0}\sum_{i=1}^n f(x_{i-1})[h(x_i)-h(x_{i-1})]\leq Mf(a).
\]
若 $f(a)=0$, 则 $f(x)\equiv0$, 结论显然成立. 若 $f(a)\neq0$, 则有
\[
  m\leq\frac{\int_a^b f(x)g(x)\mathrm dx}{f(a)}\leq M.
\]
因此存在 $\xi_2\in[a,b]$, 使得
\[
  h(\xi_2)=\int_a^{\xi_2}g(x)\mathrm dx
  =\frac{\int_a^b f(x)g(x)\mathrm dx}{f(a)}.
\]
(2) 证毕。

现在证 (3). 不妨设 $f(x)$ 单调上升, 则 $F(x)=f(x)-f(a)$ 满足 (1). 由 (1) 知, 存在 $\xi\in[a,b]$, 使得
\[
  \int_a^b F(x)g(x)\mathrm dx=F(b)\int_\xi^b g(x)\mathrm dx=[f(b)-f(a)]\int_\xi^b g(x)\mathrm dx.
\]
因此
\[
\begin{aligned}
  \int_a^b f(x)g(x)\mathrm dx
  &=f(a)\int_a^b g(x)\mathrm dx+\int_a^b F(x)g(x)\mathrm dx\\
  &=f(a)\int_a^b g(x)\mathrm dx+f(b)\int_\xi^b g(x)\mathrm dx
  -f(a)\int_\xi^b g(x)\mathrm dx\\
  &=f(a)\int_a^\xi g(x)\mathrm dx+f(b)\int_\xi^b g(x)\mathrm dx.
\end{aligned}
\]
证毕。

\noindent\textbf{例 7.5.2}\quad 设函数 $f(x)$ 在区间 $[a,b]$ 上单调上升, 证明
\[
  \int_a^b xf(x)\mathrm dx\geq\frac{a+b}{2}\int_a^b f(x)\mathrm dx.
\]

\noindent\textbf{证明}\quad 对函数 $f(x)$ 及 $g(x)=x-\dfrac{a+b}{2}$ 在区间 $[a,b]$ 上应用定积分第二中值定理知, 存在 $\xi\in[a,b]$, 使得下式成立:
\[
\begin{aligned}
  \int_a^b f(x)\left(x-\frac{a+b}{2}\right)\mathrm dx
  &=f(a)\int_a^\xi\left(x-\frac{a+b}{2}\right)\mathrm dx
  +f(b)\int_\xi^b\left(x-\frac{a+b}{2}\right)\mathrm dx\\
  &=[f(b)-f(a)]\frac{b-\xi}{2}(\xi-a)\geq0.
\end{aligned}
\]
整理即得所要结论。

\noindent\textbf{例 7.5.3}\quad 设函数
\[
  f(x)=
  \begin{cases}
    \displaystyle\int_0^x \sin\frac{1}{t}\mathrm dt, & x\neq0,\\
    0, & x=0,
  \end{cases}
\]
证明 $f(x)$ 在 $x=0$ 处可导, 并且 $f'(0)=0$。

\noindent\textbf{证明}\quad 对任给的 $x>0$, 由于 $\sin\dfrac{1}{t}$ 在区间 $[0,x]$ 上可积, 从而有
\[
  \int_0^x \sin\frac{1}{t}\mathrm dt=\lim_{\delta\to0+0}\int_\delta^x \sin\frac{1}{t}\mathrm dt.
\]
对于 $\forall 0<\delta<x$, 在上面等式右端中的定积分作积分变换 $t=\dfrac{1}{u}$, 得
\[
  \int_\delta^x \sin\frac{1}{t}\mathrm dt
  =\int_{\frac{1}{x}}^{\frac{1}{\delta}}\frac{\sin u}{u^2}\mathrm du.
\]
对上面等式右端中的定积分应用定积分第二中值定理知, 存在 $\xi_\delta\in\left[\dfrac{1}{x},\dfrac{1}{\delta}\right]$, 使得
\[
  \int_{\frac{1}{x}}^{\frac{1}{\delta}}\frac{\sin u}{u^2}\mathrm du
  =x^2\int_{\frac{1}{x}}^{\xi_\delta}\sin u\mathrm du.
\]
由于 $\left|\int_{\frac{1}{x}}^{\xi_\delta}\sin u\mathrm du\right|\leq2$, 因此, 对任意的 $\delta>0$, 有
\[
  \left|\int_\delta^x \sin\frac{1}{t}\mathrm dt\right|\leq2x^2.
\]
于是对于任意给定的 $x>0$, 令 $\delta\to0+0$, 得
\[
  \left|\int_0^x \sin\frac{1}{t}\mathrm dt\right|\leq2x^2.
\]
注意到 $\int_0^x\sin\dfrac{1}{t}\mathrm dt$ 是偶函数, 从而上式对 $x<0$ 也成立. 因此有
\[
  0\leq\lim_{x\to0}\left|\frac{\int_0^x\sin\frac{1}{t}\mathrm dt}{x}\right|
  \leq\lim_{x\to0}\frac{2x^2}{|x|}=0.
\]
从上式我们知, $f(x)$ 在 $x=0$ 处可导, 并且
\[
  f'(0)=\lim_{x\to0}\frac{\int_0^x\sin\frac{1}{t}\mathrm dt}{x}=0.
\]

\section{定积分在几何学中的应用}

\subsection{直角坐标系下平面图形的面积}

根据定积分的几何意义可以直接求出一些简单平面图形的面积. 设平面图形 $S$ 的边界是由直线 $x=a,x=b$ 及连续函数 $y=f(x),y=g(x)(x\in[a,b])$ 的图像所组成 (其中 $g(x)\leq f(x),x\in[a,b]$, 见图 7.6.1), 我们称这种形状的平面图形 $S$ 为 X 型区域. 对于这种区域的面积 $A$, 由定积分的几何意义立即得到
\[
  A=\int_a^b [f(x)-g(x)]\mathrm dx.
\]

若平面图形 $S$ 的边界是由直线 $y=c,y=d$ 及连续函数 $x=\varphi(y),x=\psi(y),y\in[c,d]$ 的图像所组成 (其中 $\varphi(y)\leq\psi(y),y\in[c,d]$, 见图 7.6.2), 则我们称这种平面图形 $S$ 为 Y 型区域. 对于这种区域的面积 $A$, 由定积分的几何意义, 我们有以下的计算公式
\[
  A=\int_c^d [\psi(y)-\varphi(y)]\mathrm dy.
\]

因此当一个平面图形 $S$ 可以分成有限个互不相交的 X 型或 Y 型区域时, 我们就可以求出 $S$ 的面积。

\noindent\textbf{例 7.6.1}\quad 求由直线 $y=2,y=x$ 及曲线 $xy=1$ 所围平面图形 $S$ 的面积 $A$。

\noindent\textbf{解}\quad 从图 7.6.3 可以看出, $S$ 是 Y 型区域, 因此有
\[
  A=\int_1^2\left(y-\frac{1}{y}\right)\mathrm dy
  =\left.\left(\frac{y^2}{2}-\ln y\right)\right|_1^2
  =\frac{3}{2}-\ln2.
\]

我们也可通过添加辅助线 $x=1$ 将图形 $S$ 分割成两个 X 型区域来进行计算:
\[
\begin{aligned}
  A&=\int_{\frac12}^1\left(2-\frac{1}{x}\right)\mathrm dx+\int_1^2(2-x)\mathrm dx\\
  &=\left.(2x-\ln x)\right|_{\frac12}^1
  +\left.\left(2x-\frac{x^2}{2}\right)\right|_1^2\\
  &=\frac{3}{2}-\ln2.
\end{aligned}
\]

\noindent\textbf{例 7.6.2(Young 不等式)}\quad 设 $y=f(x)$ 是区间 $[0,+\infty)$ 上严格上升的连续函数, 且满足 $f(0)=0$, 证明: 对任何的 $a>0,b>0$, 有
\begin{equation}
  ab\leq\int_0^a f(x)\mathrm dx+\int_0^b f^{-1}(y)\mathrm dy,
  \tag{7.6.1}
\end{equation}
并且上述不等式中等号成立的充分必要条件是 $b=f(a)$。

\noindent\textbf{证明}\quad 首先我们注意到, 若 $y=\varphi(x)$ 是区间 $[\alpha,\beta]$ 上严格上升的非负连续函数, 则由定积分的几何意义有 (见图 7.6.4)
\begin{equation}
  \beta\varphi(\beta)-\alpha\varphi(\alpha)
  =\int_\alpha^\beta \varphi(x)\mathrm dx
  +\int_{\varphi(\alpha)}^{\varphi(\beta)}\varphi^{-1}(y)\mathrm dy.
  \tag{7.6.2}
\end{equation}
现在我们回到 $f(x)$. 先设 $0<b\leq f(a)$, 由于 $f(x)$ 的严格上升性, 我们有以下不等式
\[
  [a-f^{-1}(b)]b\leq\int_{f^{-1}(b)}^a f(x)\mathrm dx.
\]
整理得
\begin{equation}
  ab-\int_0^a f(x)\mathrm dx
  \leq bf^{-1}(b)-\int_0^{f^{-1}(b)}f(x)\mathrm dx.
  \tag{7.6.3}
\end{equation}
在式 (7.6.2) 中令 $\alpha=0,\beta=b,\varphi(x)=f^{-1}(x)$, 则有
\[
  bf^{-1}(b)-\int_0^{f^{-1}(b)}f(x)\mathrm dx
  =\int_0^b f^{-1}(y)\mathrm dy,
\]
将上述等式代入 (7.6.3) 即得所证。

当 $0<f(a)\leq b$ 时, 我们可以从 $y=f^{-1}(x)$ 出发, 重复以上论证即可证明所要的结论。

显然, 当 $b=f(a)$ 时, 上面的每一不等式均成立等式, 从而 (7.6.1) 中成立等式. 反之, 若 (7.6.1) 式的等号成立, 则由 $f(x)$ 的严格上升性知必有 $b=f(a)$。

\subsection{参数方程表示的曲线所围平面图形的面积}

设平面图形 $S$ 的边界曲线 $\gamma$ 由参数方程
\[
  \begin{cases}
    x=x(t),\\
    y=y(t)
  \end{cases}
  \quad(\alpha\leq t\leq\beta)
\]
给出, 其中 $x(\alpha)=x(\beta),y(\alpha)=y(\beta)$, 曲线 $\gamma$ 除了端点重合外再无自交点 (如图 7.6.5). 假定 $x'(t),y'(t)$ 在区间 $[\alpha,\beta]$ 上存在且连续. 为了叙述简便起见, 我们称这种曲线为 $C^1$ 曲线. 再假定 $t$ 从 $\alpha$ 连续变化到 $\beta$ 时, 如果一个人沿着 $\gamma$ 行走, 图形 $S$ 总在他的左边. 这种规定给出了 $\gamma$ 的一个正定向. 以后我们总假定所求面积的图形是由满足上述条件的曲线所围。

现在假设曲线 $\gamma$:
\[
  \begin{cases}
    x=x(t),\\
    y=y(t)
  \end{cases}
  \quad(\alpha\leq t\leq\beta)
\]
围成一个 X 型区域, 如图 7.6.5 所示, 我们希望利用直角坐标求面积的公式来导出参数方程求面积的公式. 注意到当 $t$ 从 $\alpha$ 严格上升到 $t_1$ 时, $x$ 从 $x(\alpha)$ 严格下降到 $x(t_1)$, 因此 $x(t)$ 在区间 $[\alpha,t_1]$ 上存在反函数 $t=t^{-1}(x)$. 将其代入 $y=y(t)$ 得 $y=y(t^{-1}(x))\triangleq\varphi(x)(x\in[x(t_1),x(\alpha)])$. 同理, $x(t)$ 在区间 $[t_1,t_2]$ 上也存在反函数 $t^{-1}(x)$, 因此当 $t\in[t_1,t_2]$ 时, $y=y(t^{-1}(x))\triangleq\psi(x)(x\in[x(t_1),x(t_2)])$. 于是由定积分的几何意义及定积分的变量替换公式, 该 X 型区域的面积 $A$ 由以下等式表示:
\[
\begin{aligned}
  A&=\int_{x(t_1)}^{x(\alpha)}\varphi(x)\mathrm dx
  -\int_{x(t_1)}^{x(t_2)}\psi(x)\mathrm dx\\
  &=-\int_\alpha^{t_1}y(t)x'(t)\mathrm dt
  -\int_{t_1}^{t_2}y(t)x'(t)\mathrm dt
  -\int_{t_2}^{\beta}y(t)\cdot0\,\mathrm dt\\
  &=-\int_\alpha^\beta y(t)x'(t)\mathrm dt.
\end{aligned}
\]
同理, 若上述图形是一个 Y 型区域. 读者可以推出 $\gamma$ 所围区域的面积为
\[
  A=\int_\alpha^\beta x(t)y'(t)\mathrm dt.
\]
如果 $\gamma$ 所围图形可以添加几条 $C^1$ 曲线后成为若干块 X 型或 Y 型区域 (见图 7.6.6), 则 $\gamma$ 所围区域的面积仍具有形式
\begin{equation}
\begin{aligned}
  A&=\int_\alpha^\beta x(t)y'(t)\mathrm dt
  =-\int_\alpha^\beta y(t)x'(t)\mathrm dt\\
  &=\frac{1}{2}\int_\alpha^\beta [x(t)y'(t)-y(t)x'(t)]\mathrm dt.
\end{aligned}
\tag{7.6.4}
\end{equation}
这是因为 $\gamma$ 所围图形的面积是各小块图形的面积之和, 而添加的曲线必定是两块小图形的共同边界. 沿着添加的曲线的积分必须要积两次, 但由于曲线相对它所围的图形的定向相反, 因此在这些添加曲线上的积分互相抵消。

为了简便起见, 可以将 (7.6.4) 简记为
\[
  A=\int_\gamma x\mathrm dy=-\int_\gamma y\mathrm dx
  =\frac{1}{2}\int_\gamma x\mathrm dy-y\mathrm dx.
\]
若不指明 $\gamma$ 的定向, 则积分总是取 $\gamma$ 的正向。

\noindent\textbf{例 7.6.3}\quad 计算椭圆 $\dfrac{x^2}{a^2}+\dfrac{y^2}{b^2}=1$ 所围图形 $S$ 的面积 $A$(见图 7.6.7)。

\noindent\textbf{解法 1}\quad 因为 $S$ 是一个 X 型区域, 其中上、下曲线分别为 $y=\pm\dfrac{b}{a}\sqrt{a^2-x^2}$, 因此
\[
\begin{aligned}
  A&=2\int_{-a}^a\frac{b}{a}\sqrt{a^2-x^2}\mathrm dx\\
  &=\frac{2b}{a}\left.\left(\frac{x}{2}\sqrt{a^2-x^2}
  +\frac{a^2}{2}\arcsin\frac{x}{a}\right)\right|_{-a}^a
  =\pi ab.
\end{aligned}
\]

\noindent\textbf{解法 2}\quad 将椭圆方程写成参数方程
\[
  \gamma:
  \begin{cases}
    x=a\cos t,\\
    y=b\sin t,
  \end{cases}
  \quad t\in[0,2\pi],
\]
则当 $t$ 从 $0$ 连续变为 $2\pi$ 时, 该曲线为正定向. 因此
\[
\begin{aligned}
  A&=\int_\gamma x\mathrm dy
  =\int_0^{2\pi}a\cos t\cdot b\cos t\mathrm dt\\
  &=ab\int_0^{2\pi}\frac{1+\cos2t}{2}\mathrm dt
  =\pi ab.
\end{aligned}
\]

\subsection{微元法}

在以下各个小节中我们将推导一些平面或空间几何体的面积公式或体积公式. 在这里我们利用定积分的思想归纳出推导这些公式的一个普遍的方法, 即所谓的微元法。

我们考查正连续函数 $y=f(x)(x\in[a,b])$ 的定积分 $S=\int_a^b f(x)\mathrm dx$. 我们已经知道它的几何意义是由曲线 $y=f(x)(x\in[a,b])$, $x$ 轴及直线 $x=a,x=b$ 所围的面积。

对区间 $[0,S]$ 作分割 $\Delta_S:0=s_0<s_1<\cdots<s_n=S$, 则有
\begin{equation}
  \sum_{i=1}^n(s_i-s_{i-1})=\sum_{i=1}^n\Delta s_i=S=\int_a^b f(x)\mathrm dx.
  \tag{7.6.5}
\end{equation}
由于 $f(x)>0$, 因此对于 $[0,S]$ 的分割相应地有 $[a,b]$ 的分割 $\Delta:a=x_0<x_1<\cdots<x_n=b$, 使得 $\Delta s_i$ 恰好是曲线 $y=f(x)(x\in[x_{i-1},x_i])$ 与直线 $x=x_{i-1},x=x_i$ 及 $x$ 轴所围的面积. 容易看出当 $\lambda_S=\max_{1\leq i\leq n}\{\Delta s_i\}\to0$ 时, 有 $\lambda(\Delta)=\max_{1\leq i\leq n}\{\Delta x_i\}\to0$. 因此, $f(x)$ 的黎曼和
\begin{equation}
  \sum_{i=1}^n f(\xi_i)\Delta x_i\to\sum_{i=1}^n\Delta s_i=S=\int_a^b f(x)\mathrm dx.
  \tag{7.6.6}
\end{equation}
区间 $[0,S]$ 的长度和将它分成一些小区间再加起来是一样的. 现在我们换一种比较高的观点来看此问题. 任取 $[s,s+\Delta s]\subset[0,S]$, 当 $s$ 是自变量时, $\Delta s=\mathrm ds$. 因此将小区间的长度加起来, 也就是将 $\mathrm ds$ 在 $[0,S]$ 上积分起来, 即 $S=\int_0^S\mathrm ds$。

另外, 我们从 (7.6.6) 可以知道, $[s,s+\Delta s]$ 对应有一个 $[x,x+\Delta x]$, 使得 $\Delta s=f(\xi)\Delta x$, 其中 $\xi\in[x,x+\Delta x]$. 当 $\Delta s$ 很小时, $\Delta x$ 也很小. 现在我们将 $x$ 看成自变量, 则有 $\Delta x=\mathrm dx$. 若将 $\xi$ 取成 $x$, 则由于 $\Delta x\to0$ 时有 $\Delta s=f(x)\Delta x+o(\Delta x)$. 由此我们有 $\mathrm ds=f(x)\mathrm dx$. 因此得到
\begin{equation}
  S=\int_0^S\mathrm ds=\int_a^b f(x)\mathrm dx.
  \tag{7.6.7}
\end{equation}

将上面的过程总结一下. 我们希望求 $S$, 即曲线 $y=f(x)(x\in[a,b])$ 与直线 $x=a,x=b$ 及 $y=0$ 所围的曲边梯形的面积. 假定我们不知道求该面积的公式, 现在来找出该公式. 为此我们任取一小曲边梯形: 由过 $x,x+\Delta x$ 的两条平行于 $y$ 轴的直线、直线 $y=0$ 与曲线 $y=f(x)$ 所构成的小曲边梯形 (如图 7.6.8). 同时得到 $[0,S]$ 上的一个小区间 $[s,s+\Delta s]$, 使得 $\Delta s$ 等于该小曲边梯形的面积, 即有
\[
  \Delta s=f(\xi)\Delta x\approx f(x)\Delta x=f(x)\mathrm dx,
\]
其中 $\xi$ 是 $[x,x+\Delta x]$ 中的某一点. 如果 $f(x)$ 连续, 则有 $\mathrm ds=f(x)\mathrm dx$. 因此我们得到了 (7.6.7)。

去掉上述例子的具体背景, 我们可以归纳出以下求某些量的一般方法. 设所求的量是 $S$, 它与区间 $[a,b]$ 有关, 当区间 $[a,b]$ 给定后, $S$ 就是一个确定的量, 而且该量对该区间具有可加性, 即对于 $[a,b]$ 的任意分割 $\Delta:a=x_0<x_1<\cdots<x_n=b$, 若记 $[x_{i-1},x_i]$ 对应的部分量为 $\Delta s_i(i=1,2,\cdots,n)$, 则 $S=\sum_{i=1}^n\Delta s_i$. 为求 $S$, 我们可以任取小区间 $[x,x+\mathrm dx]\subset[a,b]$, 若 $[x,x+\mathrm dx]$ 所对应的部分量 $\Delta s=f(\xi)\mathrm dx(\xi\in[x,x+\mathrm dx])$, 其中 $f(x)$ 是 $[a,b]$ 上的一个连续函数, 则我们即有 $\mathrm ds=f(x)\mathrm dx$, 并且 $S=\int_a^b f(x)\mathrm dx$. 这个求 $S$ 的过程也称微元法, 其中 $\mathrm ds=f(x)\mathrm dx$ 称为量 $S$ 的微元. 微元法在后续章节中也多次用到。

\subsection{极坐标方程表示的曲线所围平面图形的面积}

在本小节中, 我们设曲线 $\gamma$ 由 $r=r(\theta)(\alpha\leq\theta\leq\beta)$ 所表示, 其中 $r(\theta)$ 是 $\theta$ 的连续函数. 现在我们来求曲线 $\gamma$ 与射线 $\theta=\alpha,\theta=\beta$ 所围平面图形 (见图 7.6.9) 的面积 $A$。

我们用微元法来推导其面积公式. 任取 $[\theta,\theta+\mathrm d\theta]\subset[\alpha,\beta]$. 因此由曲线 $r=r(\theta)$ 及向径 $\theta_1=\theta,\theta_2=\theta+\mathrm d\theta$ 所围的图形近似于一个三角形, 其高近似为 $r(\theta)$, 而底边长近似为 $r(\theta)\mathrm d\theta$, 因此其面积微元为
\[
  \mathrm dA=\frac{1}{2}r^2(\theta)\mathrm d\theta.
\]
从 $\alpha$ 到 $\beta$ 将其面积微元积分得
\[
  A=\int_\alpha^\beta\frac{1}{2}r^2(\theta)\mathrm d\theta.
\]

\noindent\textbf{例 7.6.4}\quad 求心脏线 $r=a(1+\cos\theta)(a>0,0\leq\theta\leq2\pi)$ 所围平面图形 $S$ 的面积 $A$。

\noindent\textbf{解}\quad
\[
  A=\frac{1}{2}\int_0^{2\pi}a^2(1+\cos\theta)^2\mathrm d\theta=\frac{3}{2}\pi a^2.
\]

\noindent\textbf{例 7.6.5}\quad 求由圆弧
\[
  \begin{cases}
    x=\cos t,\\
    y=\sin t
  \end{cases}
  \quad\left(0\leq t\leq\frac{3}{4}\pi\right)
\]
及 $x$ 轴和直线 $x=-\dfrac{\sqrt2}{2}$ 所围平面图形的面积 (见图 7.6.10)。

\noindent\textbf{解法 1}\quad 我们用参数方程所围平面图形的面积公式来求, 注意到在线段 $BC$ 上有 $\mathrm dx=0$, 而在线段 $CD$ 上有 $y=0$, 因此该图形的面积为
\[
\begin{aligned}
  A&=-\int_0^{\frac34\pi}y\mathrm dx
  =-\int_0^{\frac34\pi}\sin t(-\sin t)\mathrm dt\\
  &=\int_0^{\frac34\pi}\frac{1-\cos2t}{2}\mathrm dt
  =\left.\left(\frac{t}{2}-\frac{\sin2t}{4}\right)\right|_0^{\frac34\pi}\\
  &=\frac{3}{8}\pi+\frac{1}{4}.
\end{aligned}
\]

\noindent\textbf{解法 2}\quad 现在我们用极坐标表示的曲线所围平面图形的面积公式. 参数方程表示的圆弧方程为 $r=1\left(0\leq\theta\leq\dfrac{3}{4}\pi\right)$. 所围图形面积 $A$ 应为扇形 $BOD$ 的面积加上三角形 $OBC$ 的面积:
\[
  A=\int_0^{\frac34\pi}\frac{1}{2}\cdot1^2\mathrm d\theta
  +\frac{1}{2}\left(\frac{\sqrt2}{2}\right)^2
  =\frac{3}{8}\pi+\frac{1}{4}.
\]

\subsection{平行截面面积为已知的立体的体积}

利用定积分, 我们还可以求出空间一些规则立体图形的体积. 设空间某立体 $\Omega$ 介于平面 $x=a$ 和 $x=b$ 之间, 假定过 $x\in[a,b]$ 并且与 $x$ 轴垂直的平面截该立体 $\Omega$ 的截面的面积为已知函数 $A(x)$(见图 7.6.11), 进一步假定 $A(x)$ 是区间 $[a,b]$ 上的连续函数. 现在我们来求 $\Omega$ 的体积 $V$。

利用微元法, 任取 $[x,x+\mathrm dx]\subset[a,b]$, $\Omega$ 夹在过点 $x=x$ 和 $x=x+\mathrm dx$ 并且与 $x$ 轴垂直的两个平面之间的部分近似于一个柱体, 因此体积微元 $\mathrm dV=A(x)\mathrm dx$. 将体积微元从 $a$ 到 $b$ 积分得
\[
  V=\int_a^b A(x)\mathrm dx.
\]
特别地, 若该立体 $\Omega$ 的边界曲面是由平面 $x=a,x=b$ 以及 $Oxy$ 平面内一条曲线 $y=f(x)\geq0(x\in[a,b])$ 绕 $x$ 轴旋转而成, 则称这个立体为旋转体. 这种旋转体 $\Omega$ 与过 $x\in[a,b]$ 并且与 $x$ 轴垂直的平面相截的面积为 $A(x)=\pi f^2(x)$. 因此 $\Omega$ 的体积 $V$ 由下面公式给出:
\[
  V=\pi\int_a^b f^2(x)\mathrm dx.
\]

\noindent\textbf{例 7.6.6}\quad 求椭球体 $\dfrac{x^2}{a^2}+\dfrac{y^2}{b^2}+\dfrac{z^2}{c^2}\leq1(a,b,c>0)$ 的体积 $V$。

\noindent\textbf{解}\quad 对于每个 $x\in[-a,a]$, 过点 $x$ 且与 $x$ 轴垂直的平面截该椭球体的截面为椭圆
\[
  \frac{y^2}{b^2\left(1-\frac{x^2}{a^2}\right)}
  +\frac{z^2}{c^2\left(1-\frac{x^2}{a^2}\right)}=1
\]
所围的平面图形, 由例 7.6.2 知该图形的面积为 $\pi bc\left(1-\dfrac{x^2}{a^2}\right)$, 因此
\[
  V=\int_{-a}^a \pi bc\left(1-\frac{x^2}{a^2}\right)\mathrm dx
  =\left.\pi bc\left(x-\frac{x^3}{3a^2}\right)\right|_{-a}^a
  =\frac{4}{3}\pi abc.
\]

\noindent\textbf{例 7.6.7}\quad 求星形线 $x^{\frac23}+y^{\frac23}=a^{\frac23}(a>0)$ 绕 $x$ 轴旋转一周所成的立体体积 $V$。

\noindent\textbf{解}\quad 由旋转体的体积公式, 我们有
\[
  V=\pi\int_{-a}^a y^2\mathrm dx
  =\pi\int_{-a}^a (a^{\frac23}-x^{\frac23})^3\mathrm dx
  =\frac{32}{105}\pi a^3.
\]

\subsection{曲线的弧长}

设平面曲线 $\Gamma$ 由参数方程
\[
  \begin{cases}
    x=x(t),\\
    y=y(t)
  \end{cases}
  \quad(\alpha\leq t\leq\beta)
\]
给出. 若 $x(t)$ 与 $y(t)$ 是区间 $[\alpha,\beta]$ 上的连续函数, 则我们称 $\Gamma$ 是一条连续曲线 (简称曲线). 我们下面涉及的曲线, 除特别说明外, 均指连续曲线。

下面我们来讨论曲线的弧长问题. 为此我们将面临以下三个问题:

(1) 如何来定义曲线的弧长?

(2) 是不是曲线都有弧长?

(3) 对有弧长的曲线怎样来求出它的弧长?

首先来讨论问题 (1). 设平面曲线 $\Gamma$ 的起点为 $A$, 终点为 $B$. 由于我们目前能够精确求出线段长度, 因此希望用线段的长度去逼近曲线的长度. 为此我们在 $\Gamma$ 上从 $A$ 到 $B$ 依次取 $n+1$ 个点 $A=M_0,M_1,M_2,\cdots,M_{n-1},M_n=B$. 这 $n+1$ 个点将 $\Gamma$ 分成了 $n$ 段小弧段. 记 $L_i(i=1,2,\cdots,n)$ 为以 $M_{i-1},M_i$ 两点为端点的线段的长度. 如果分点无限增加, 且当 $\lambda=\max_{1\leq i\leq n}\{L_i\}\to0$ 时, $\sum_{i=1}^n L_i$ 有极限 $L$, 则称 $\Gamma$ 为可求长曲线, 并且定义 $\Gamma$ 的长度为 $L$。

现在我们来讨论问题 (2): 是不是曲线都可以求长度呢? 为此我们考虑由下述参数方程表示的曲线 $\Gamma$: 当 $t\in(0,1]$ 时, $x=t,y=t\sin\dfrac{1}{t}$; 当 $t=0$ 时, $x=y=0$。

我们取 $t'_n=\dfrac{1}{2n\pi+\frac12\pi}$, $t''_n=\dfrac{1}{2n\pi-\frac12\pi}(n=1,2,\cdots,N)$, 设 $\Gamma$ 上对应 $t'_n$ 的点为 $M'_n$, 而对应 $t''_n$ 的点为 $M''_n$, 则我们有
\[
  \sum_{n=1}^N L_n
  >\sum_{n=1}^N\left|t'_n\sin\frac{1}{t'_n}-t''_n\sin\frac{1}{t''_n}\right|
  =\sum_{n=1}^N(t'_n+t''_n)>\frac{1}{2\pi}\sum_{n=1}^N\frac{1}{n},
\]
其中 $L_n(n=1,2,\cdots,N)$ 是连接 $M'_n$ 和 $M''_n$ 的线段的长度. 因为当分点个数 $N\to\infty$ 时, 不等式右端将趋于 $+\infty$, 从而 $\Gamma$ 是不可求长的曲线。

从上面的例子可知, 并非所有的曲线都是可求长的, 我们必须对曲线加上适当条件才能使得它是可求长的。

下面我们将讨论问题 (3). 如果我们要对一般的可求长曲线来求出它的长度, 除了利用弧长的定义外, 似乎很难找到其他的求弧长公式. 为了利用微积分的工具来讨论曲线的求长问题, 我们必须对可求长曲线加上适当的光滑条件, 因此在下面主要讨论光滑曲线的求长问题. 设 $\Gamma$:
\[
  \begin{cases}
    x=x(t),\\
    y=y(t)
  \end{cases}
  \quad(\alpha\leq t\leq\beta)
\]
是一条曲线, 如果它满足 $x'(t),y'(t)$ 在区间 $[\alpha,\beta]$ 上存在且连续, 并且 $[x'(t)]^2+[y'(t)]^2\neq0$, 则称 $\Gamma$ 是一条光滑曲线。

我们下面还是用微元法来推导这类曲线的弧长公式. 任取 $[t,t+\mathrm dt]\subset[\alpha,\beta]$, 得到曲线上两个点 $M_t(x(t),y(t)),M_{t+\mathrm dt}(x(t+\mathrm dt),y(t+\mathrm dt))$, 则弦长 $\overline{M_tM_{t+\mathrm dt}}$ 有以下计算公式:
\[
\begin{aligned}
  \overline{M_tM_{t+\mathrm dt}}
  &=\sqrt{[x(t+\mathrm dt)-x(t)]^2+[y(t+\mathrm dt)-y(t)]^2}\\
  &=\sqrt{[x'(\xi_1)]^2+[y'(\xi_2)]^2}\mathrm dt\\
  &\approx\sqrt{[x'(t)]^2+[y'(t)]^2}\mathrm dt,
\end{aligned}
\]
其中 $\xi_1,\xi_2$ 在 $t$ 与 $t+\mathrm dt$ 之间. 因此, 当 $\mathrm dt$ 很小时, 两点之间的弧长 $\Delta L$ 可以由弦长近似表示:
\[
  \Delta L\approx\sqrt{[x'(t)]^2+[y'(t)]^2}\mathrm dt.
\]
这样我们得到弧长微元如下:
\[
  \mathrm dL=\sqrt{[x'(t)]^2+[y'(t)]^2}\mathrm dt.
\]
因此 $\Gamma$ 的弧长公式为
\[
  L=\int_\alpha^\beta\sqrt{[x'(t)]^2+[y'(t)]^2}\mathrm dt.
\]
特别地, 当曲线 $\Gamma$ 为函数 $y=f(x)(x\in[a,b])$ 的图像时, 若 $f(x)$ 在 $[a,b]$ 上具有连续导数, 则有
\[
  L=\int_a^b\sqrt{1+[f'(x)]^2}\mathrm dx.
\]
如果曲线 $\Gamma$ 由极坐标 $r=r(\theta)(\alpha\leq\theta\leq\beta)$ 方程给出, 其中 $r'(\theta)$ 在 $\alpha\leq\theta\leq\beta$ 上存在并且连续, 则我们可将其写成参数式
\[
  \begin{cases}
    x=r(\theta)\cos\theta,\\
    y=r(\theta)\sin\theta
  \end{cases}
  \quad(\alpha\leq\theta\leq\beta),
\]
从而弧长微元为
\[
\begin{aligned}
  \mathrm dL
  &=\sqrt{[r'(\theta)\cos\theta-r(\theta)\sin\theta]^2
  +[r'(\theta)\sin\theta+r(\theta)\cos\theta]^2}\mathrm d\theta\\
  &=\sqrt{[r(\theta)]^2+[r'(\theta)]^2}\mathrm d\theta.
\end{aligned}
\]
因此 $\Gamma$ 的长度 $L$ 有下面的计算公式:
\[
  L=\int_\alpha^\beta\sqrt{[r(\theta)]^2+[r'(\theta)]^2}\mathrm d\theta.
\]
如果 $\Gamma$ 是空间曲线, $\Gamma$ 由参数方程
\[
  \begin{cases}
    x=x(t),\\
    y=y(t),\\
    z=z(t)
  \end{cases}
  \quad(\alpha\leq t\leq\beta)
\]
给出, 并且 $x'(t),y'(t),z'(t)$ 在区间 $[\alpha,\beta]$ 上连续且不同时为零, 则 $\Gamma$ 的弧长 $L$ 有下面的计算公式:
\[
  L=\int_\alpha^\beta\sqrt{[x'(t)]^2+[y'(t)]^2+[z'(t)]^2}\mathrm dt.
\]

\noindent\textbf{例 7.6.8}\quad 求曲线 $\Gamma:y=x^2(x\in[0,1])$ 的弧长 $L$。

\noindent\textbf{解}\quad
\[
\begin{aligned}
  L&=\int_0^1\sqrt{1+[(x^2)']^2}\mathrm dx
  =\int_0^1\sqrt{1+4x^2}\mathrm dx\\
  &=\left.\left[x\sqrt{x^2+\frac{1}{4}}
  +\frac{1}{4}\ln\left(x+\sqrt{x^2+\frac{1}{4}}\right)\right]\right|_0^1\\
  &=\frac{\sqrt5}{2}+\frac{1}{4}\ln\left(1+\frac{\sqrt5}{2}\right)-\frac{1}{4}\ln\frac{1}{2}.
\end{aligned}
\]

\noindent\textbf{例 7.6.9}\quad 求双纽线 $r^2=2a^2\cos2\theta(a>0)$ 从 $\theta=0$ 到 $\theta=\dfrac{\pi}{6}$ 之间的弧的长度 $L$。

\noindent\textbf{解}\quad 对 $r^2=2a^2\cos2\theta$ 两边关于 $\theta$ 求导得 $2rr'=-4a^2\sin2\theta$, 因此
\[
  r'(\theta)=\frac{-2a^2\sin2\theta}{r},
\]
且
\[
\begin{aligned}
  \sqrt{r^2(\theta)+[r'(\theta)]^2}
  &=\sqrt{r^2+\frac{(-2a^2\sin2\theta)^2}{r^2}}
  =\sqrt{\frac{r^4+4a^4\sin^22\theta}{r^2}}\\
  &=\frac{\sqrt2a}{\sqrt{\cos2\theta}}
  =\frac{\sqrt2a}{\sqrt{1-2\sin^2\theta}}.
\end{aligned}
\]
最后得到
\[
  L=\int_0^{\frac{\pi}{6}}\frac{\sqrt2a}{\sqrt{1-2\sin^2\theta}}\mathrm d\theta.
\]
此积分称为椭圆积分, 由于被积函数没有初等原函数, 因此需要用其他的方法算出其值。

\subsection{旋转体的侧面积}

下面我们来讨论旋转体的侧面积问题. 设 $\Gamma$ 为一光滑曲线, 其参数方程为
\[
  \begin{cases}
    x=x(t),\\
    y=y(t)
  \end{cases}
  \quad(\alpha\leq t\leq\beta).
\]
再设当 $t\in[\alpha,\beta]$ 时, $y=y(t)\geq0$, 则曲线 $\Gamma$ 绕 $x$ 轴旋转一周后与 $x=x(\alpha)$ 和 $x=x(\beta)$ 两个平面围成一个旋转体 $\Omega$. 现在我们用微元法来推导 $\Omega$ 的侧面积 $S$ 的计算公式。

任取 $[t,t+\mathrm dt]\subset[\alpha,\beta]$, 曲线 $\Gamma$ 在 $[t,t+\mathrm dt]$ 的部分绕 $x$ 轴旋转所形成的立体可近似看做一个圆台 (图 7.6.12), 其侧面积为
\[
  \Delta S\approx\pi[y(t)+y(t+\mathrm dt)]\overline{M(t)M(t+\mathrm dt)},
\]
其中 $M(t)=(x(t),y(t)),M(t+\mathrm dt)=(x(t+\mathrm dt),y(t+\mathrm dt))$. 当 $\mathrm dt$ 很小时, $\overline{M(t)M(t+\mathrm dt)}\approx\mathrm dL$, 其中 $\mathrm dL=\sqrt{[x'(t)]^2+[y'(t)]^2}\mathrm dt$ 是 $\Gamma$ 的弧长微元, 而 $y(t+\mathrm dt)\approx y(t)$, 因此我们得到面积微元为
\[
  \mathrm dS=2\pi y(t)\sqrt{[x'(t)]^2+[y'(t)]^2}\mathrm dt.
\]
对 $t$ 从 $\alpha$ 到 $\beta$ 积分即得所求侧面积:
\[
  S=2\pi\int_\alpha^\beta y(t)\sqrt{[x'(t)]^2+[y'(t)]^2}\mathrm dt.
\]
类似地, 我们可以给出曲线绕 $y$ 轴旋转一周所得的旋转体的侧面积公式。

对于在其他坐标下的曲线绕坐标轴旋转一周所得立体的侧面积公式, 请读者自己给出。

\noindent\textbf{例 7.6.10}\quad 求圆周 $x^2+y^2=R^2$ 上 $x$ 介于 $[a,a+h]$(其中 $-R\leq a<a+h\leq R$) 的部分绕 $x$ 轴旋转一周所形成的球台的体积及侧面积。

\noindent\textbf{解}\quad 该球台可由曲线
\[
  y=f(x)=\sqrt{R^2-x^2},\quad x\in[a,a+h]
\]
绕 $x$ 轴旋转一周而成, 因此其体积
\[
\begin{aligned}
  V&=\pi\int_a^{a+h}(R^2-x^2)\mathrm dx
  =\left.\pi\left(R^2x-\frac{x^3}{3}\right)\right|_a^{a+h}\\
  &=\pi\left[R^2h-\left(a^2h+ah^2+\frac{h^3}{3}\right)\right].
\end{aligned}
\]
而该球台的侧面积为
\[
\begin{aligned}
  S&=2\pi\int_a^{a+h}\sqrt{R^2-x^2}\sqrt{1+[f'(x)]^2}\mathrm dx\\
  &=2\pi R\int_a^{a+h}\frac{\sqrt{R^2-x^2}}{\sqrt{R^2-x^2}}\mathrm dx
  =2\pi hR.
\end{aligned}
\]

\noindent\textbf{注}\quad 在上述结果中, 令 $a=-R,h=2R$, 则我们得到半径为 $R$ 的球的体积为 $\dfrac{4\pi R^3}{3}$, 球面面积为 $4\pi R^2$. 另外, 我们还可看出球台的体积不仅依赖于球台的高, 而且还依赖于 $a$ 的选取. 但它的侧面积则为球面上大圆周长与球台的高的乘积, 而与 $a$ 的选取无关。

\section{定积分在物理学中的应用}

定积分在物理学中有着广泛的应用. 在本节中, 我们主要讨论一些规则物体的质心、转动惯量的计算问题。

关于这些物理量的计算, 读者首先应该知道在有限个质点的情形下相应物理量的计算公式, 然后再根据定积分的定义来推导出其他形状物体的相应物理量的计算公式。

先来考虑质心坐标问题. 首先来复习一下物理学的有关概念. 设平面上有 $n$ 个质点, 它们在平面内的位置为 $A_1(x_1,y_1),A_2(x_2,y_2),\cdots,A_n(x_n,y_n)$, 在点 $A_i$ 处的质点的质量为 $m_i(i=1,2,\cdots,n)$. 由物理学中的定义, 这 $n$ 个质点关于 $x$ 轴和 $y$ 轴的静力矩分别为
\[
  M_x=\sum_{i=1}^n m_i y_i\quad\text{和}\quad M_y=\sum_{i=1}^n m_i x_i.
\]
记 $M=\sum_{i=1}^n m_i$ 为这 $n$ 个质点的总质量, 则我们有以下的静力矩定律:

在上述条件下, 存在点 $A(\overline x,\overline y)$, 满足
\[
  M_x=\sum_{i=1}^n m_i y_i=M\overline y,\qquad
  M_y=\sum_{i=1}^n m_i x_i=M\overline x.
\]
因此我们可得到以下结论: 这 $n$ 个质点关于 $x$ 轴 ($y$ 轴) 的静力矩之和相当于一个位于在点 $(\overline x,\overline y)$ 处质量为 $M$ 的质点关于 $x$ 轴 ($y$ 轴) 的静力矩. 因此 $A(\overline x,\overline y)$ 可以认为是这 $n$ 个质点的质心, 其坐标由以下公式给出:
\[
  \overline x=\frac{M_y}{M}
  =\frac{\sum_{i=1}^n m_i x_i}{\sum_{i=1}^n m_i},\qquad
  \overline y=\frac{M_x}{M}
  =\frac{\sum_{i=1}^n m_i y_i}{\sum_{i=1}^n m_i}.
\]
在物理学中, 我们还分别称
\[
  I_x=\sum_{i=1}^n m_i y_i^2\quad\text{与}\quad I_y=\sum_{i=1}^n m_i x_i^2
\]
为这 $n$ 个质点关于 $x$ 轴与 $y$ 轴的转动惯量。

现在我们来求平面内一物质曲线的质心和关于坐标轴的转动惯量. 所谓物质曲线是一个理想的模型, 即平面 (或空间) 有一条光滑曲线 $\Gamma$, 它由参数方程 $x=x(t),y=y(t)(\alpha\leq t\leq\beta)$ 给出, 同时在 $\Gamma$ 上有一个质量分布函数, 它由线密度 $\rho(t)(\alpha\leq t\leq\beta)$ 给出, 其中 $\rho(t)\in C[\alpha,\beta]$。

由微元法, 任取 $[t,t+\mathrm dt]\subset[a,b]$, 则曲线弧长微元
\[
  \mathrm dL=\sqrt{[x'(t)]^2+[y'(t)]^2}\mathrm dt.
\]
因此得到质量微元
\[
  \mathrm dM=\rho(t)\mathrm dL=\rho(t)\sqrt{[x'(t)]^2+[y'(t)]^2}\mathrm dt,
\]
从而整个曲线的质量为
\[
  M=\int_\alpha^\beta \rho(t)\sqrt{[x'(t)]^2+[y'(t)]^2}\mathrm dt.
\]
另外, 由质点的静力矩的定义易知, 该物质曲线关于 $x$ 轴及 $y$ 轴的静力矩微元
\[
  \mathrm dM_x=\mathrm dM\cdot y(t),\qquad
  \mathrm dM_y=\mathrm dM\cdot x(t),
\]
从而 $\Gamma$ 关于 $x$ 轴和 $y$ 轴的静力矩分别为
\[
  M_x=\int_\alpha^\beta \rho(t)y(t)\sqrt{[x'(t)]^2+[y'(t)]^2}\mathrm dt,
\]
\[
  M_y=\int_\alpha^\beta \rho(t)x(t)\sqrt{[x'(t)]^2+[y'(t)]^2}\mathrm dt.
\]
由此我们可以求出 $\Gamma$ 的质心坐标 $(\overline x,\overline y)$:
\[
  \overline x=\frac{M_y}{M}
  =\frac{\int_\alpha^\beta \rho(t)x(t)\sqrt{[x'(t)]^2+[y'(t)]^2}\mathrm dt}
  {\int_\alpha^\beta \rho(t)\sqrt{[x'(t)]^2+[y'(t)]^2}\mathrm dt},
\]
\[
  \overline y=\frac{M_x}{M}
  =\frac{\int_\alpha^\beta \rho(t)y(t)\sqrt{[x'(t)]^2+[y'(t)]^2}\mathrm dt}
  {\int_\alpha^\beta \rho(t)\sqrt{[x'(t)]^2+[y'(t)]^2}\mathrm dt}.
\]
若在上面公式中, 特别地令 $\rho(t)\equiv1$, 即曲线 $\Gamma$ 具有均匀的质量分布, 则得到质心坐标
\[
  \overline x=\frac{\int_\alpha^\beta x(t)\sqrt{[x'(t)]^2+[y'(t)]^2}\mathrm dt}{L},
\]
\[
  \overline y=\frac{\int_\alpha^\beta y(t)\sqrt{[x'(t)]^2+[y'(t)]^2}\mathrm dt}{L},
\]
其中 $L$ 是曲线 $\Gamma$ 的长度。

特别地, 我们有
\[
  2\pi\overline yL
  =2\pi\int_\alpha^\beta y(t)\sqrt{[x'(t)]^2+[y'(t)]^2}\mathrm dt=S,
\]
其中 $S$ 是曲线 $\Gamma$ 绕 $x$ 轴旋转一周得到的旋转体的侧面积. 因此上面的等式告诉我们, 一条曲线绕 $x$ 轴旋转一周所成旋转体的侧面积等于该曲线长度与质心绕 $x$ 轴旋转一周的周长的乘积. 这个结果也称为是古鲁金 (Guldin) 第一定理。

对于物质曲线 $\Gamma$ 关于 $x$ 轴和 $y$ 轴的转动惯量, 我们用同样的方法可得到如下公式:
\[
  I_x=\int_\alpha^\beta y^2(t)\rho(t)\sqrt{[x'(t)]^2+[y'(t)]^2}\mathrm dt,
\]
\[
  I_y=\int_\alpha^\beta x^2(t)\rho(t)\sqrt{[x'(t)]^2+[y'(t)]^2}\mathrm dt.
\]

\noindent\textbf{例 7.7.1}\quad 求圆周
\[
  x^2+(y-a)^2=R^2\quad(a>R>0)
\]
绕 $x$ 轴旋转所成的轮胎体的表面积 (见图 7.7.1)。

\noindent\textbf{解}\quad 将该圆周看成线密度函数 $\rho\equiv1$ 的质量曲线. 显然, 它的质心坐标为 $(0,a)$. 注意到圆的周长为 $L=2\pi R$, 因此由古鲁金第一定理得其表面积为
\[
  S=L\cdot2\pi\overline y=2\pi R\cdot2\pi a=4\pi^2Ra.
\]

现在我们来考虑平面图形的质心问题. 设平面图形 $S$ 由曲线 $y=y_1(x),y=y_2(x)(y_1(x)\leq y_2(x),a\leq x\leq b)$ 及直线 $x=a,x=b$ 所围成, 并设该平面图形 $S$ 的质量为均匀分布, 且假定其面密度函数为 $1$。

下面我们仍由微元法来求其质心坐标. 任取 $[x,x+\mathrm dx]\subset[a,b]$, 则该图形 $S$ 落在过点 $x$ 和 $x+\mathrm dx$ 并且与 $x$ 轴垂直的直线之间的部分可近似看成一个小矩形, 从而质量微元为 $\mathrm dM=[y_2(x)-y_1(x)]\mathrm dx$, 质心坐标近似为
\[
  \left(x,\frac{1}{2}[y_1(x)+y_2(x)]\right)
\]
(见图 7.7.2). 因此得到 $S$ 关于 $x$ 轴和 $y$ 轴的静力矩微元分别为
\[
\begin{aligned}
  \mathrm dM_x
  &=\frac{1}{2}[y_2(x)+y_1(x)][y_2(x)-y_1(x)]\mathrm dx\\
  &=\frac{1}{2}[y_2^2(x)-y_1^2(x)]\mathrm dx,
\end{aligned}
\]
\[
  \mathrm dM_y=x[y_2(x)-y_1(x)]\mathrm dx.
\]
对 $x$ 从 $a$ 到 $b$ 积分即得到整个图形 $S$ 关于 $x$ 轴和 $y$ 轴的静力矩分别为
\[
  M_x=\frac{1}{2}\int_a^b [y_2^2(x)-y_1^2(x)]\mathrm dx,
\]
\[
  M_y=\int_a^b x[y_2(x)-y_1(x)]\mathrm dx.
\]
注意到 $S$ 的总质量为
\[
  M=\int_a^b [y_2(x)-y_1(x)]\mathrm dx,
\]
因此该图形 $S$ 的质心坐标为
\[
  \overline x=\frac{M_y}{M}
  =\frac{\int_a^b x[y_2(x)-y_1(x)]\mathrm dx}
  {\int_a^b [y_2(x)-y_1(x)]\mathrm dx},
\]
\[
  \overline y=\frac{M_x}{M}
  =\frac{\dfrac{1}{2}\int_a^b [y_2^2(x)-y_1^2(x)]\mathrm dx}
  {\int_a^b [y_2(x)-y_1(x)]\mathrm dx}.
\]

由上面关于 $\overline y$ 的计算公式可得
\[
\begin{aligned}
  2\pi\overline yS
  &=2\pi\overline y\int_a^b [y_2(x)-y_1(x)]\mathrm dx\\
  &=\pi\int_a^b [y_2^2(x)-y_1^2(x)]\mathrm dx=V.
\end{aligned}
\]
其中 $S$ 是该图形的面积, $V$ 是该图形绕 $x$ 轴旋转一周所得旋转体的体积. 上述公式表明一个平面图形绕 $x$ 轴旋转一周得到的旋转体的体积等于该平面图形的面积与以质心到 $x$ 轴距离为半径的圆周长的乘积. 这个结果也称为古鲁金第二定理。

\noindent\textbf{例 7.7.2}\quad 求半径为 $R>0$ 的半圆盘的质心。

\noindent\textbf{解}\quad 如图 7.7.3 建立坐标系. 设半圆盘的质心坐标为 $(\overline x,\overline y)$, 由对称性显然有 $\overline x=0$. 下面求 $\overline y$. 注意到该图形绕 $x$ 轴旋转一周得到一个球体, 利用古鲁金第二定理有
\[
  2\pi\overline y\cdot\frac{\pi}{2}R^2=\frac{4}{3}\pi R^3.
\]
解上述方程得 $\overline y=\dfrac{4}{3\pi}R$, 因此圆盘质心为 $\left(0,\dfrac{4}{3\pi}R\right)$。

下面我们举一个例子来讨论极坐标方程表示的物质曲线转动惯量的计算。

\noindent\textbf{例 7.7.3}\quad 设平面质量曲线为 $r=r(\theta)(\alpha\leq\theta\leq\beta)$, 线密度为 $\rho\equiv1$, 求该曲线关于极轴的转动惯量。

\noindent\textbf{解}\quad 任取 $[\theta,\theta+\mathrm d\theta]\subset(\alpha,\beta)$, 则由转动惯量的定义, 曲线上位于 $[\theta,\theta+\mathrm d\theta]$ 上的弧段质量微元为 $\mathrm dM=\sqrt{r^2(\theta)+r'^2(\theta)}\mathrm d\theta$. 这小段质量弧到极轴的距离近似为 $r(\theta)\sin\theta$ (见图 7.7.4), 它关于极轴的转动惯量, 即转动惯量微元为
\[
  \mathrm dI=r^2(\theta)\sin^2\theta\,\mathrm dM,
\]
因此所求曲线关于极轴的转动惯量为
\[
  I=\int_\alpha^\beta r^2(\theta)\sin^2\theta
  \sqrt{r^2(\theta)+r'^2(\theta)}\mathrm d\theta.
\]

\section*{习题七}

1. 从定义出发计算定积分 $\int_0^T(\alpha+\beta t)\mathrm dt$, 这里 $\alpha,\beta$ 是常数。

\end{document}
